\documentclass[10pt,a4paper]{article}

% ----------------- Paquetes -------------------%
\usepackage[T1]{fontenc}
\usepackage[utf8]{inputenc}
\usepackage[a4paper,margin=2.5cm]{geometry}
\usepackage{mathptmx}             
\usepackage{changepage} 
\usepackage{setspace}
\usepackage[spanish]{babel}
\usepackage[labelfont=bf,labelsep=space]{caption}
\usepackage{amssymb}
\usepackage{amsmath}
\usepackage{ebgaramond}
\usepackage{placeins}
\usepackage{etoolbox}
\usepackage{setspace}
\usepackage{parskip} 
\usepackage{tcolorbox}
\usepackage{needspace}
\usepackage{xparse}
\usepackage{ocgx2}
\usepackage{hyperref}
\usepackage{hypcap}
\usepackage{soul}\sethlcolor{yellow}% resaltado amarillo: \hl{texto} (Panel TCEE, ⌃H)



%%%%%%%%%%%%%%%%%%%%%%%%%%%%%%%%%%%%%%%%%%%%%%%%%%%%%%%%%%%%%%%%
% --- Fija primero tus valores globales "buenos" ---
\setstretch{1.2}                 % Interlineado global
\setlength{\parskip}{0.7\baselineskip}
\setlength{\parindent}{0pt}
\newlength{\GlobalParskip}
\setlength{\GlobalParskip}{\parskip}

\newcounter{eqnum}[subsection]
\renewcommand{\theeqnum}{\arabic{eqnum}}
\newcounter{alphaitem}
\newcounter{numitem}

%% ------------------- Recuadros----------------------%%
\newcommand{\puntosclave}[1]{%
  \begin{tcolorbox}[
    colframe=black,
    title={Puntos clave},
    before upper={\setlength{\parindent}{0.5cm}\setlength{\parskip}{0.5\baselineskip}}
  ]
  #1
  \end{tcolorbox}
}

\newcommand{\modificaciones}[1]{
  \begin{tcolorbox}[
    colback=white,
    colframe=red,
    title={Modificaciones a realizar},
    before upper={\setlength{\parindent}{0.5cm}\setlength{\parskip}{0.5\baselineskip}}
  ]
  #1
  \end{tcolorbox}
}

%% ------------------- Preguntas TEST----------------------%%
% Opciones: radio (single-choice)
\NewDocumentCommand{\OptRadio}{m m m}{%
  \noindent\CheckBox[radio,name=#1,value=#2,width=1.2ex,height=1.2ex]{}%
  \;\textbf{#2.}~#3\par
}

% Opciones: checkbox (multi-choice)
\NewDocumentCommand{\OptCheck}{m m}{%
  \noindent\CheckBox[name=#1,width=1.2ex,height=1.2ex,bordercolor={0 0 0}]{}%
  \;#2\par
}

% Pregunta single-choice (radio)
% Args: 1=id, 2=enunciado, 3=opciones, 4=solución+justificación, 5=imagen (o {}), 6=origen
\NewDocumentCommand{\PreguntaTestSingle}{m m m m m m}{%
  \par\medskip
  \noindent\textbf{#1.}~#2\par\smallskip

    % Imagen (si no es {})
    \IfBlankTF{#5}{}{%
      \begin{center}
        \includegraphics[width=0.75\linewidth]{#5}
      \end{center}
    }

    \begin{Form}
      #3
    \end{Form}

    \par\smallskip
    \switchocg{sol-#1}{\fbox{Mostrar / ocultar solución}}%
    \hfill{\footnotesize #6}\par\smallskip

    \begin{ocg}{Solución}{sol-#1}{0}
      \noindent #4
    \end{ocg}
  \par\medskip
}

% Pregunta multi-choice (checkboxes)
\NewDocumentCommand{\PreguntaTestMulti}{m m m m m m}{%
  \par\medskip
  \noindent\textbf{#1.}~#2\par\smallskip

    % Imagen (si no es {})
    \IfBlankTF{#5}{}{%
      \begin{center}
        \includegraphics[width=0.75\linewidth]{#5}
      \end{center}
    }
    \begin{Form}
      #3
    \end{Form}

    \par\smallskip
    \switchocg{sol-#1}{\fbox{Mostrar / ocultar solución}}%
    \hfill{\footnotesize #6}\par\smallskip

    \begin{ocg}{Solución}{sol-#1}{0}
      \noindent #4
    \end{ocg}
  \par\medskip
}

%% ------------------- Listas----------------------%%
    % --- Lista alfabética ---
    \newenvironment{la}
      {\begin{list}{\alph{alphaitem})}{%
          \usecounter{alphaitem}%
          \setlength{\leftmargin}{1cm}%
          \setlength{\parskip}{3pt}% 
          \setlength{\itemsep}{0pt}% ← espacio entre ítems
          \setlength{\parsep}{0pt}%  ← párrafos dentro de un ítem
      }}
      {\end{list}}
      
    % --- Lista numérica ---
    \newenvironment{lnum}
      {\begin{list}{\arabic{numitem}.}{%
          \usecounter{numitem}%
          \setlength{\leftmargin}{1cm}%
          \setlength{\parskip}{3pt}% 
          \setlength{\itemsep}{0pt}% ← espacio entre ítems
          \setlength{\parsep}{0pt}%  ← párrafos dentro de un ítem
      }}
      {\end{list}}

%% ------------------- Preguntas Test ----------------------%%
      \usepackage{xparse} % si ya lo tienes, no lo repitas

    \NewDocumentEnvironment{test}{m m o}{%
      \begin{tcolorbox}[
        colback=white,
        colframe=black!15,
        boxrule=0.6pt,
        arc=2mm,
        left=2mm,right=2mm,top=2mm,bottom=2mm
      ]
      \centering
      \includegraphics[width=\linewidth]{#1}
      \end{tcolorbox}
    
      \vspace{2mm}
    
      % Caja SOLO para texto (SÍ partible y mantiene márgenes al partir)
      \begin{tcolorbox}[
        enhanced,
        breakable,
        colback=black!3,
        colframe=black!25,
        boxrule=0.6pt,
        arc=2mm,
        left=2mm,right=2mm,top=1.5mm,bottom=1.5mm
      ]
      \textbf{Respuesta:} \textbf{#2}%
      \IfValueT{#3}{\par\smallskip \textbf{Justificación:} #3}
      \end{tcolorbox}
    }{}
      
% ------------------- Imágenes----------------------%
    \graphicspath{{Img/}}% Rutas por defecto 
    % --- Imagen adaptativa ---
    % Registros de dimensión definidos una sola vez
    \newdimen\imgcap
    \newdimen\imgmin
    \newdimen\imgmax
    \newdimen\imgremain
    \newdimen\imgh
    \newdimen\imgneed
    
    \newcommand{\imagenfit}[3]{%
      \addvspace{10pt}% separación antes
      \par\begingroup
        % Parámetros de composición (asignaciones, no \newdimen)
        \imgcap=1\baselineskip            % alto estimado de título+fuente
        \imgmin=0.15\textheight
        \imgmax=0.15\textheight
        \imgremain=\pagegoal
        \ifdim\pagegoal=\maxdimen \imgremain=0pt\fi
        \advance\imgremain by -\pagetotal
        \advance\imgremain by -\imgcap
        % Decisión de altura destino
        \ifdim\imgremain<\imgmin
          \newpage
          \imgh=0.20\textheight
        \else
          \imgh=\imgremain
          \ifdim\imgh>\imgmax \imgh=\imgmax \fi
        \fi
        % Reservar espacio para evitar cortes del bloque completo
        \imgneed=\imgh \advance\imgneed by \imgcap
        \Needspace{\imgneed}
        % Cuerpo
        \noindent\begin{minipage}{\textwidth}
          \centering
          \captionof{figure}{\textemdash\ #2}\par\vspace{0.3em}% título arriba
          \includegraphics[width=\linewidth,height=\imgh,keepaspectratio]{\detokenize{#1}}\par
          \vspace{0.3em}\small\textit{Fuente: }#3% fuente abajo
        \end{minipage}%
      \endgroup
      \par\addvspace{10pt}%
    }

%------------ Bloque de RECUADRO ESPECIAL -------------%
    \newenvironment{specialblock}[1]{%
      \begin{quote} % crea sangría a izquierda y derecha
      \small % texto más pequeño
      \noindent\textbf{#1}\\[-0.3em] % título en negrita
      \rule{\linewidth}{0.4pt}\\[0.6em]}{
      \end{quote}}
      
%-------------------------- Portada -----------------------%
\newcommand{\oldtitle}[1]{\def\OldTitle{#1}}
\newcommand{\changerationale}[1]{\def\ChangeWhy{#1}}

\newcommand{\changebox}{%
\begin{tcolorbox}[title={Historial de cambios del tema}]
\textbf{Título anterior:} \OldTitle\par
\textbf{Motivo del cambio:} \ChangeWhy
\end{tcolorbox}
}

%-------------------------- Author Font -----------------------%
    \newcommand{\authorfont}[1]{{\fontfamily{EBGaramond-TLF}\selectfont #1}}

%-------------------------- Anexos -----------------------%
    \makeatletter
    \newcounter{anexo}
    \renewcommand{\theanexo}{\Roman{anexo}}
    \newcommand{\anexo}[1]{%
      \clearpage
      \phantomsection
      \refstepcounter{anexo}%
      \noindent\textbf{Anexo \theanexo.- }#1\par\medskip
      \addcontentsline{stc}{section}{Anexo \theanexo.- #1}%
    }
    \makeatother

    %-------------------------- Lista de figuras (personal) -----------------------%
\makeatletter
\newcommand{\MyListOfFigures}{%
  \clearpage
  \phantomsection
  {\centering\bfseries\Large Índice de figuras\par}%
  \medskip
  % Entrada en nuestro índice de contenido (stc)
  \addcontentsline{stc}{section}{Índice de figuras}%
  % Recuperación del listado (sin cabecera propia)
  \begingroup
    \parindent=0pt\parskip=2pt
    \@starttoc{lof}%
  \endgroup
}
\makeatother

%-------------------------- Bibliografía -----------------------%
\makeatletter
% Contador para las referencias
\newcounter{myref}
% Cabecera de la bibliografía, entra en el índice 'stc'
\newcommand{\MyBibliography}{%
  \clearpage
  \phantomsection
  {\centering\bfseries\Large Bibliografía y referencias\par}%
  \medskip
  \addcontentsline{stc}{section}{Bibliografía y referencias}%
  \setcounter{myref}{0}%
}
\newcommand{\MyBibItem}[4]{%
  \refstepcounter{myref}%
  \noindent[\themyref]\ #1.\ \emph{#2}. #3. #4\par
}
\makeatother

%--------- Establecer cuatro niveles de índice --------%
    \setcounter{tocdepth}{4}   
    \setcounter{secnumdepth}{4}% numera hasta \paragraph

%%%%%%%%%%%%%%%%%%%%%%%%%%%%%%%%

\makeatletter

    %--------- Interlineado Global --------%
    \edef\GlobalBaseLineStretch{\baselinestretch}
    
    %--------- Espaciado de seccinones --------%
    \renewcommand\section{\@startsection{section}
      {1}{\z@}
      {2.5ex \@plus 1ex \@minus .2ex} % espacio antes
      {1.5ex \@plus .2ex}             % espacio después
      {\normalfont\Large\bfseries}    % formato del título
    }
    \renewcommand\subsection{\@startsection{subsection}{2}{\z@}
      {3ex \@plus 0.8ex \@minus .2ex}
      {1.5ex \@plus .2ex}
      {\normalfont\large\bfseries}
    }
    \renewcommand\subsubsection{\@startsection{subsubsection}{3}{\z@}
      {3ex \@plus 0.5ex \@minus .2ex}
      {1ex \@plus .2ex}
      {\normalfont\normalsize\bfseries}
    }
    \renewcommand\paragraph{\@startsection{paragraph}{4}{\z@}%
    {-2ex \@plus -0.5ex \@minus -.2ex}% espacio antes
    {0.8ex \@plus .2ex}% espacio después
    {\normalfont\normalsize\itshape}}% ← cursiva en lugar de negrita

    %--------- Ecuaciones --------%   
    % Variables globales para almacenar títulos y notas
    \gdef\leq@current@section{}%
    \gdef\leq@current@subsection{}%
    \gdef\leq@last@printed@section{}%
    \gdef\leq@last@printed@subsection{}%
    
    % Comando para añadir nota (invisible en el cuerpo)
    \newcommand{\eqnote}[1]{%
      \addcontentsline{leq}{leqnote}{#1}%
    }
    
    % Comando para ecuaciones
    \newcommand{\eqblock}[2]{%
      % Verificar si necesitamos imprimir el título de sección
      \ifx\leq@current@section\leq@last@printed@section\else
        \addcontentsline{leq}{leqsection}{\leq@current@section}%
        \global\let\leq@last@printed@section\leq@current@section
        \gdef\leq@last@printed@subsection{}% Reset subsección al cambiar sección
      \fi
      % Verificar si necesitamos imprimir el título de subsección
      \ifx\leq@current@subsection\leq@last@printed@subsection\else
        \ifx\leq@current@subsection\@empty\else
          \addcontentsline{leq}{leqsubsection}{\leq@current@subsection}%
          \global\let\leq@last@printed@subsection\leq@current@subsection
        \fi
      \fi
      % Ecuación
      \refstepcounter{eqnum}%
      \begin{equation*}
        #1\tag{E.\theeqnum}
      \end{equation*}%
      \addcontentsline{leq}{leqt}{[E.\theeqnum] -- #2}%
      \addcontentsline{leq}{leqm}{#1}%
    }
    
    %%% ----- Índice de ecuaciones ----------%%%
    % Tamaño para []
    \newcommand{\leqIndexFont}{\normalsize}
    \newcommand{\leqTitleFont}{\tiny}
    \newcommand{\leqNoteFont}{\tiny}
    % Cabecera de sección 
    \newcommand*\l@leqsection[2]{%
      \addvspace{25pt}% Mayor separación antes de nueva sección
      \noindent\leqIndexFont
      \makebox[\linewidth]{\rule{.38\linewidth}{0.4pt}\ \ \textbf{  \MakeUppercase{#1}}\ \ \rule{.38\linewidth}{0.4pt}}%
      \par\vspace{-10pt}
    }
    % Cabecera de subsección 
    \newcommand*\l@leqsubsection[2]{%
      \par\addvspace{15pt}%
      \noindent\leqIndexFont #1
      \par\vspace{-7pt}
    }
    % Título de ecuación 
    \newcommand*\l@leqt[2]{%
    \par\addvspace{10pt}%
      \par\noindent\leqTitleFont #1\par
    }
    % Miniatura de ecuación
    \newcommand*\l@leqm[2]{%
      \vspace{0.3\baselineskip}%
      \par\scriptsize\noindent\hspace*{1.5em}\(#1\)\par%
    }
    % Nota de ecuación 
    \newcommand*\l@leqnote[2]{%
      \vspace{0.2\baselineskip}%
      \par\noindent\leqNoteFont\hspace*{1.5em}\textbf{Nota.--} #1\par\vspace{0.5\baselineskip}%
    }
    % Generación del índice
    \newcommand{\mylistofequations}{%
      \clearpage
      \phantomsection
      % Título visual de la lista (ajústalo a tu gusto):
      {\centering\bfseries\Large Índice de ecuaciones\par}
      % Entrada en nuestro índice 'stc':
      \addcontentsline{stc}{section}{Índice de ecuaciones}%
      % Llamada a tu lista real (usa el nombre exacto de tu macro):
      \@starttoc{leq}%
    }
    
    %--------- Captura de títulos de sección --------%
    \let\leq@original@section\section
    \let\leq@original@subsection\subsection
    % Flag para saber si estamos en una sección asterisco
    \newif\ifLeqStarSection
    \LeqStarSectionfalse
    
    % Redefinir \section CON CONTROL DE ASTERISCO
    \renewcommand{\section}{%
      \@ifstar{\leq@section@star}{\@dblarg\leq@section@nostar}%
    }
    \def\leq@section@star#1{%
      \global\LeqStarSectiontrue
      \gdef\leq@current@section{#1}%
      \gdef\leq@current@subsection{}%
      \leq@original@section*{#1}%
      \global\LeqStarSectionfalse
    }
    \def\leq@section@nostar[#1]#2{%
      \global\LeqStarSectionfalse
      \gdef\leq@current@section{#2}%
      \gdef\leq@current@subsection{}%
      \leq@original@section[#1]{#2}%
    }
    % Redefinir \subsection
    \renewcommand{\subsection}{%
      \@ifstar{\leq@subsection@star}{\@dblarg\leq@subsection@nostar}%
    }
    \def\leq@subsection@star#1{%
      \global\LeqStarSectiontrue
      \gdef\leq@current@subsection{#1}%
      \leq@original@subsection*{#1}%
      \global\LeqStarSectionfalse
    }
    \def\leq@subsection@nostar[#1]#2{%
      \global\LeqStarSectionfalse
      \gdef\leq@current@subsection{#2}%
      \leq@original@subsection[#1]{#2}%
    }

    % ----- Restauración de valores Globales de estilo ----------%
    % Flags
    \newif\ifInFootnote
    \InFootnotefalse
    \pretocmd{\@footnotetext}{\InFootnotetrue}{}{}
    \apptocmd{\@footnotetext}{\InFootnotefalse}{}{}
    % Profundidad de listas (soporta anidadas)
    \newcount\MyListDepth \MyListDepth=0
    \AtBeginEnvironment{la}{\advance\MyListDepth by 1}
    \AfterEndEnvironment{la}{\advance\MyListDepth by -1}
    \AtBeginEnvironment{lnum}{\advance\MyListDepth by 1}
    \AfterEndEnvironment{lnum}{\advance\MyListDepth by -1}
    \AtBeginEnvironment{itemize}{\advance\MyListDepth by 1}
    \AfterEndEnvironment{itemize}{\advance\MyListDepth by -1}
    \AtBeginEnvironment{enumerate}{\advance\MyListDepth by 1}
    \AfterEndEnvironment{enumerate}{\advance\MyListDepth by -1}
    % Restauración diferida: hazlo al INICIO del próximo párrafo real
    \newif\if@pendingRestore
    \@pendingRestorefalse
    \newcommand{\RequestRestore}{\global\@pendingRestoretrue}
    \everypar{%
      \if@pendingRestore
        \global\@pendingRestorefalse
        \ifInFootnote\else
          \ifnum\MyListDepth=0
            \setlength{\parskip}{\GlobalParskip}%
            \setstretch{\GlobalBaseLineStretch}%
          \fi
        \fi
      \fi
    }
    % Al final de bloques:
    \AfterEndEnvironment{equation}{\RequestRestore}
    \AfterEndEnvironment{equation*}{\RequestRestore}
    \AfterEndEnvironment{la}{\RequestRestore}
    \AfterEndEnvironment{lnum}{\RequestRestore}
    % Al final de macros:
    \apptocmd{\eqblock}{\RequestRestore}{}{}
    \apptocmd{\imagenfit}{\RequestRestore}{}{}
    
\makeatother

% ===================== ToC propio con 4 niveles (único bloque) =====================
\makeatletter

% --- Escritores: añadimos una línea al .stc en cada cabecera numerada ---
% Guardamos las versiones actuales para llamar al comportamiento normal
\let\MyTOC@orig@section\section
\let\MyTOC@orig@subsection\subsection
\let\MyTOC@orig@subsubsection\subsubsection
\let\MyTOC@orig@paragraph\paragraph

% SECTION
\renewcommand{\section}{%
  \@ifstar{\MyTOC@section@star}{\@dblarg\MyTOC@section@nostar}%
}
\def\MyTOC@section@star#1{%
  \MyTOC@orig@section*{#1}% sin entrada al .stc
}
\def\MyTOC@section@nostar[#1]#2{%
  \MyTOC@orig@section[{#1}]{#2}% comportamiento normal (escribe al .toc)
  % Escribimos también nuestra línea al .stc (texto con \numberline)
  \addcontentsline{stc}{section}{\protect\numberline{\thesection}#2}%
}

% SUBSECTION
\renewcommand{\subsection}{%
  \@ifstar{\MyTOC@subsection@star}{\@dblarg\MyTOC@subsection@nostar}%
}
\def\MyTOC@subsection@star#1{%
  \MyTOC@orig@subsection*{#1}%
}
\def\MyTOC@subsection@nostar[#1]#2{%
  \MyTOC@orig@subsection[{#1}]{#2}%
  \addcontentsline{stc}{subsection}{\protect\numberline{\thesubsection}#2}%
}

% SUBSUBSECTION
\renewcommand{\subsubsection}{%
  \@ifstar{\MyTOC@subsubsection@star}{\@dblarg\MyTOC@subsubsection@nostar}%
}
\def\MyTOC@subsubsection@star#1{%
  \MyTOC@orig@subsubsection*{#1}%
}
\def\MyTOC@subsubsection@nostar[#1]#2{%
  \MyTOC@orig@subsubsection[{#1}]{#2}%
  \addcontentsline{stc}{subsubsection}{\protect\numberline{\thesubsubsection}#2}%
}

% PARAGRAPH
\renewcommand{\paragraph}{%
  \@ifstar{\MyTOC@paragraph@star}{\@dblarg\MyTOC@paragraph@nostar}%
}
\def\MyTOC@paragraph@star#1{%
  \MyTOC@orig@paragraph*{#1}%
}
\def\MyTOC@paragraph@nostar[#1]#2{%
  \MyTOC@orig@paragraph[{#1}]{#2}%
  % Si no numeras los \paragraph, cambia \theparagraph por texto plano sin \numberline
  \addcontentsline{stc}{paragraph}{\protect\numberline{\theparagraph}#2}%
}

% --- Impresor del índice propio (lee el .stc) ---
\newcommand{\MyTOC}{%
  \par\bigskip
  {\centering\bfseries\Large Índice de contenido\par}%
  \medskip
  \begingroup
    \parindent=0pt\parskip=2pt
    \@starttoc{stc}%
  \endgroup
}

\makeatother
% ===================== fin ToC propio =====================


%%%%%%%%%%%%%%


% --- Datos del tema ---
\title{Tema 3.A.20 -- Regulación y Poder de mercado\\
\large }
\author{Marcelo Ortega}
\date{Febrero 2026}
\newcommand{\TemaCode}{3A12}

\oldtitle{Tema 3.A.8. La dualidad en la teoría de la demanda del consumidor y sus aplicaciones}
\changerationale{
    Se explicitan los contenidos necesarios del tema para guiar al opositor y garantizar un mínimo común. Se añade una referencia a aplicaciones empíricas de sistemas de demanda, siguiendo la línea iniciada por Angus Deaton.
}

\begin{document}


\maketitle
\changebox

\newpage

% --- Página de resumen y modificaciones ---
\puntosclave{ %% Escribir puntos clave Apuntes CECO
    El tratamiento de la definición de mercado relevante debería ser rápido 
    
    Igualmente, la Teoría tradicional de la regulación (donde cabría destacar, principalmente, los precios de RAMSEY-BOITEAUX, por su influencia en desarrollos posteriores). 
    
    Dedicar más tiempo a la regulación mediante incentivos, que introduce más complejidad pero también resulta más cercano a los problemas a los que se enfrenta la regulación en la práctica. 
    
    Las principales ideas del Modelo LAFFONT-TIROLE (1986) deberían de exponerse; sin necesidad de desarrollar éste completamente. Las relaciones verticales y precios de acceso quizás deberían tener un papel menos relevante (tratarse con menos detalle). 
    
    Los mercados de varias caras (multi-sided markets) deberían tratarse con más detalle y, sobretodo, actualizarse una vez entre en aplicación el DMA (Digital Markets Act, Reglamento de Mercados Digitales, UE) \textcolor{red}{\textbf{[OJO! 2023]}}
    }
\vspace{1cm}

\modificaciones{
    El apartado de bundling and unbundling debe recuperarse, extenderse y eliminarse el Tema 3.A.17: Está incluido como mecanismo de discriminación de precios y debe ser incluido como técnica de utilización del poder de mercado
    
    Interesante incluir algo de esto: \href{https://regulationbodyofknowledge.org}{Body of Knowledge on Infrastructure Regulation}
    }

\newpage
\MyTOC
\newpage

\mylistofequations
\clearpage

\MyListOfFigures
\clearpage


\pagenumbering{arabic}
\setcounter{page}{1}


\section{Introducción}\label{intro}


\subsection{Contextualización}\label{context}


\subsubsection{Relevancia}\label{importance}
Una rama de investigación activa en los últimos años se plantea si el poder de mercado es un asunto microeconómico o si, por el contrario, también tiene consecuencias macroeconómicas. Como se estudia en los temas de macroeconomía, el poder de mercado es una característica básica de los modelos actuales de equilibrio general dinámico y estocástico (DSGE), con la introducción de la formulación a la Dixit y Stiglitz. El Fondo Monetario Internacional, en el World Economic Outlook (2019) dedicó un capítulo a las consecuencias macroeconómicas del creciente poder de mercado.

Las tres conclusiones principales de un amplio análisis de la situación de las empresas en distintos países son las siguientes:
\begin{lnum}
    \item El poder de mercado ha aumentado moderadamente en las economías avanzadas, como indica el aumento de los \textit{mark-ups} de las empresas por encima de los costes marginales en cerca de un 8\% desde el 2000. Pero no así en las economías emergentes;
    \item El aumento ha sido bastante generalizado en todas las economías e industrias avanzadas, pero dentro de ellas se ha concentrado en una pequeña fracción de empresas dinámicas, más productivas e innovadoras; y,
    \item Aunque las implicaciones macroeconómicas generales han sido modestas hasta ahora, un mayor aumento del poder de mercado de estas empresas podría aumentar estructuralmente la inflación, debilitar la inversión, desalentar la innovación, reducir la proporción de los salarios en la renta nacional y dificultar la estabilización del output por el Banco Central
\end{lnum}

Además, varios estudios muestran que la alta concentración empresarial podría tener efectos perniciosos en el mercado laboral, al favorecer las características de monopsonio de las empresas dominantes, influyendo sobre la cantidad y el precio de equilibrio (a la baja). Dentro del mercado laboral, también habría que estar al debate sobre los \textit{gig-workers} o trabajadores de plataformas, y las oportunidades y riesgos que introduce este nuevo tipo de relaciones laborales;\footnote{
\href{https://www.economist.com/finance-and-economics/2018/10/27/economists-think-antitrust-policy-should-pay-more-attention-to-workers}{\textit{Economists think antitrust policy should pay more attention to workers} The Economist (2018)}
}



\subsubsection{Historia}\label{history}



\subsubsection{Problemática}\label{problem} 




\subsection{Estructura de la exposición}\label{structure}




\clearpage
\section{Regulación: ¿Dónde, por qué y cuándo?}
\puntosclave{
    Un mercado relevante es el conjunto de productos y zonas geográficas que ejercen una presión competitiva mutua.
    
	La técnica más habitual para definir un mercado (de producto y geográfico) es el test SSNIP.
    
	La definición del mercado relevante se emplea como paso previo a la medición del poder de mercado, es decir, la capacidad de una empresa para aumentar sus precios por encima de su coste marginal.
    
	La intervención pública en presencia de poder de mercado se suele materializar en la política de competencia y la regulación.
    
	La política regulatoria puede justificarse desde un punto de vista normativo (o de interés público) o positivo (partiendo de la teoría de la regulación económica).
}

Cuando hablamos de poder de mercado, nos centramos en desgranar las causas y consecuencias de éste. Esto hace que a menudo pase desapercibido una cuestión clave: ¿qué mercado? Y, en particular, ¿qué es el mercado? Esto no es solo crucial porque constituye la condición sine qua non para hablar de \textit{poder}, sino también porque nos permite centrar rigurosamente el debate en las caraterísticas concretas del \textbf{mercado} y, como tal, las causas y consecuencias de este poder: no es lo mismo hablar del poder que ejerce el badulaque de una pequeña aldea sobre los habitantes de ésta, en cuanto único proveedor (mercado geográfico), que hablar del poder que puede ostentar mercadona en el sector de supermercados (mercado de producto). Evidentemente, las causas que han conducido a dicho poder son muy diferentes y, por la misma razón, no lo son menos las consecuencias.

\subsection{Mercado relevante: SSNIP}
Un mercado relevante es el conjunto de productos y/o zonas geográficas que ejercen una presión competitiva mutua.\footnote{%
    Para una definición en profundidad, véase la Comunicación de la Comisión relativa a la definición de mercado de referencia a efectos de la normativa comunitaria en materia de competencia.
}

Lo habitual para definir ambos mercados es partir del test SSNIP (\textit{Small but Significant Nontransitory Increase in Price}), también llamado test del monopolista hipotético. El test SSNIP parte de la idea de que los productos que pertenecen a un mismo mercado no se ven influidos en sus decisiones de precios por los productos de otros mercados.\footnote{%
    Por ejemplo, no es probable que el precio de las baterías suponga una restricción para el precio del agua embotellada, por lo que no formarían parte del mismo mercado. En cambio, el precio del agua con gas podría tener una influencia, algo habría que comprobar mediante el test SSNIP.
}

\subsubsection{Mercado de producto}
Si quisieramos definir el mercado de producto, el ejercicio consistiría en partir de una cesta reducida productos y valorar si un monopolista hipotético tiene o no incentivos a aumentar su precio de forma no transitoria y significativa (se suele tomar como referencia un aumento del 5-10\%):\footnote{%
    Al partir de un grupo de productos y considerar qué otros productos podrían formar parte del mismo mercado de referencia, se suele comenzar con aquellos productos que los consumidores consideran como sustitutivos, lo que se denominamos sustituibilidad de la demanda. No obstante, también hay que considerar la sustituibilidad de la oferta, es decir, aquellos productores que actualmente no operan en el mercado pero que serían capaces (debido a sus conocimientos, sus activos productivos, etc.) de empezar a producir los productos de referencia en un periodo relativamente corto si se produjera un aumento de precio.
}
\begin{la}
    \item En caso afirmativo, se consideraría que los productos pertenecen a un mismo mercado; y,
    \item Por el contrario, si una parte sustancial de su demanda se desplazaría a otros productos, concluiríamos que el mercado de producto inicialmente propuesto no constituye un mercado relevante. Si quisieremos continuar, por tanto, habría que ampliar la cesta de productos (el mercado de producto) y repetir el test SSNIP hasta que el aumento de precios fuera provechoso.
\end{la}

Pese a sus problemas, el test SSNIP sigue siendo ampliamente utilizado como base teórica para definir los mercados relevantes.\footnote{%
    El más sonado sería la \textit{falacia del celofán}, que nos advierte del cuidado que hay que tener cuando se aplica el test fuera del ámbito del control de concentraciones. Supongamos que nos encontramos en un mercado en el que hay una única empresa instalada que se observa que no sería rentable que incrementase 5-10\% de su precio. En este caso, no podemos concluir automáticamente que la empresa no tiene poder de mercado y que es necesario ampliar la definición del mercado, ya que la empresa puede haber establecido un precio supracompetitivo que maximiza sus beneficios precisamente porque tiene poder de mercado. Por tanto, en estos casos, para aplicar el test SSNIP hay que partir del precio competitivo, lo cual puede ser complicado en la práctica.

También podríamos considerar: (i) Test de correlación de precios. La idea es que, si dos productos pertenecen al mismo mercado, sus precios deberían moverse de forma similar a lo largo del tiempo; o, Otras técnicas. Por ejemplo, el análisis de las características físicas, del uso que se le da a los productos o de las preferencias de los consumidores (pueden estudiarse mediante encuestas).
} Al basarse en un escenario hipotético hacen falta datos y herramientas que complican su aplicabilidad real, entre los que encontramos: (i) elasticidad-precio propia, uno de los datos más importantes, ya que nos informa de cómo varía la demanda de un producto ante un cambio en su precio;\footnote{%
    Sin embargo, no basta con observar cómo varían las cantidades vendidas ante cambios en los precios, ya que puede que simultáneamente hayan variado otros factores que afectan a la demanda. Por ello, puede ser necesario recurrir a una estimación econométrica de la elasticidad para tratar de aislar el resto de los factores relevantes.
} y, (ii) elasticidad-precio cruzada, que ayuda a valorar si otros productos son sustitutivos más o menos cercanos, cuanto menor sea mayor será la indicación de que pertenecen a mercados distintos.

\subsubsection{Mercado geográfico}
En caso de querer analizar el mercado desde una perspetica geográfica (de localización), el test SSNIP seguiría siendo el punto de partida: empezar por un mercado geográfico lo más estrecho posible e ir ampliándolo. 

A la hora de aplicar el test a la definición de un mercado geográfico se podrían utilizar otros métodos como el análisis de los flujos comerciales (por ejemplo, si apenas hay importaciones y exportaciones en un área geográfica, puede que se trate de un mercado relevante) o de la importancia de los costes de transporte como porcentaje del precio del producto.\footnote{%
    Por ejemplo, muchos mercados de distribución comercial (supermercados, gasolineras) suelen tener una dimensión local o incluso inferior, puesto que la mayoría de los consumidores no querrá desplazarse grandes distancias para comprar. En cambio, el mercado de construcción de aviones se define a nivel global, ya que los costes de transporte son muy bajos respecto a los precios.
}

\subsection{Poder de mercado}
El poder de mercado puede definirse como la capacidad de una empresa para aumentar sus precios por encima de su coste marginal. Desde una perspectiva teórica, se suele medir el poder de mercado mediante el Índice de Lerner: 

\eqblock{
\begin{aligned}
L_i=\left(\frac{P_x-CMg_i}{P_x}\right)
\end{aligned}
}{Definición teórica (Índice de LERNER). Poder de Mercado: Cuantificación}

Como resulta obvio, la aplicación práctica del Índice de LERNER es complicada, entre otras cosas por la dificultad de estimar el coste marginal de una empresa.  Por ello, se suele intentar medir el poder de mercado de forma indirecta mediante el análisis de las cuotas de mercado de las empresas de interés, que se pueden obtener una vez definido el mercado: 

\eqblock{
\begin{aligned}
    L_i=\left(\frac{P_x-CMg_i}{P_x}\right)=\frac{s_i}{|\varepsilon_{p,x}|}
\end{aligned}
}{Aproximación Índice de LERNER por cuotas de mercado. Poder de Mercado: Cuantificación}

\subsubsection{Implicaciones de Política Económica}
Ahora bien, ¿por qué es necesario intervenir/regular en una situación de poder de mercado? ¿cómo se justifica? 

La intervención se justifica utilizando como marco de referencia el mercado en competencia perfecta: la fijación de un precio superior al coste marginal genera una ineficiencia asignativa conocida como pérdida irrecuperable de eficiencia. De esta forma, la intervención quedaría justificado por ser el equilibrio resultante una asignación ineficiente en el sentido de PARETO.\footnote{%
    Aunque existe alguna situación particular en la que el poder de mercado se utilizase de manera tal que se consiguiese la eficiencia asignativa característica de la competencia perfecta (la discriminación de primer grado), la obviaremos por darse únicamente bajo supuestos muy específicos y demasiado restrictivos.
} 

Sin embargo, esto no nos debe parecer suficiente. Pensemos en lo siguiente: ¿no es \textit{normal} y \textit{recurrente} la aparición de beneficios extraordinarios? ¿No forma parte de la dinámica económica? ¿No eliminaría esto el incentivo empresarial? Es más, ¿qué pasa con las patentes? ¿son entonces estas herramientas perversas que pone a disposición el mismo regulador a las empresas?

\paragraph{¿Cuándo?: Metodología de la NTOI}
Resulta evidente que justificar la intervención no basta para saber \textbf{cuándo} y \textbf{cómo} se debe plantear. Necesitamos una guía de acción más concreta que nos permita saber dónde, cuándo y cómo hay que intervenir para minimizar las ineficiencias asignativas en situaciones de poder de mercado. La corriente de la NTOI nos proporciona esta herramienta. Esta Escuela da un giro profundo en la concepción del poder de mercado y centra su atención en la conducta de las empresas: la ineficiencia no está en la estructura, sino en la conducta. De una manera simplificada pero muy ilustrativa, podemos pensar en el siguiente resumen de criterios normativos para la intervención:

\vspace{1em}
\begin{center}\setlength{\tabcolsep}{2pt}
    \begin{tabular}{|p{0.25\linewidth}|p{0.15\linewidth}|p{0.25\linewidth}|p{0.25\linewidth}|}
        \hline
        \textbf{La competencia...} &  & \multicolumn{2}{c|}{\textit{¿Es deseable?}} \\
        \hline
         &  & \textbf{Sí} & \textbf{No} \\
        \hline
        \multirow{2}{*}{\textit{¿Es factible/posible?}} 
        & \textbf{Sí} & {\scriptsize Competencia} & {\scriptsize Exceso de Entrada} \\
        \cline{2-4}
        & \textbf{No} & {\scriptsize Barreras a la Entrada} & {\scriptsize Monopolio Natural} \\
        \hline
    \end{tabular}
\end{center}
\vspace{1em}

Con esta sencilla tabla, se sintetizan las causas de la ineficiencia que deben guiar al regulador en su intervención. Es más, en ella subyace la naturaleza de la intervención que se debe llevar a cabo. ¿Debemos aplicar el mismo criterio al intervenir sobre un mercado donde los incumbentes están realizando estrategias dirigidas a desincentivar posibles entrantes que sobre un mercado en el que el grado de competencia está generando externalidades negativas sin mejorar el bienestar social? Aterrizado, ¿debemos utilizar el mismo criterio al intervenir en el mercado de la pesca que sobre las grandes plataformas de \textit{streaming}? La respuesta es: no. De hecho, ya en la segunda revolución industrial se distinguían diferentes \textit{poderes de mercado} que debían considerarse de manera separada y con una paleta de instrumentos diferenciada:
\begin{la}
    \item \textbf{Políticas en defensa de la competencia}. La defensa de la competencia, o interveciones \textit{antitrust}, tienen el objetivo de incentivar la competencia en mercados donde esta es factible y socialmente deseable. Recurriendo a nuestra tabla, consistiría en instrumentos dirigidos a las situaciones donde la competencia es desable pero no factible/posible. Aunque no existe un instrumento único para la intervención, se caracteriza por regulaciones destinadas a limitar la adquisición, protección y extensión del poder de mercado estableciendo prohibiciones a ciertos tipos de conductas. El ejemplo más sonado al respecto fue el de Standard Oil, la fragmentación de Standard Oil buscaba reducir el poder de mercado que esta había adquirido habida cuenta de que la competencia era tecnológica y legalmente posible; y,
    \item \textbf{La regulación}. La otra vertiente consiste en la regulación. Los objetivas de esta vertiente son más claros y por ello consta de una paleta de instrumentos bien definida. Como este va a ser el objeto de nuestro tema, diremos simplemente que: (i) estudia las formas óptimas de intervenir en mercados donde la competencia es posible/factible pero socialmente no deseable (en exceso);\footnote{%
        El ejemplo más claro al respecto es el de la pesca. Se trata de un sector con costes fijos relativamente bajos cuyo asset principal es el kno-how del operador y las condiciones marítimas del país. Preguntémonos: ¿sería deseable una competencia cruda y total que bajará los costes del pescado a expensas de sobreexplotar las reservas marítimas del territorio? Es precisamente porque la respuesta más lógica e intuitiva es que no, por lo que existe, de hecho, una regulación a nivel europeo al respecto.
    } y, (ii) mercados en los que tecnológicamente no es posible (ni deseable, consecuentemente) tener competencia. Estos últimos mercados se caracterizan por la presencia de lo que se conoce como \textbf{monopolioi natural}.
\end{la}

En definitiva, la cuestión es que la complejidad de las estructuras de mercado exige que cada una sea tratada de forma particular, sin perjuicio de que las	regularidades presentes en términos conceptuales nos permitan establecer un marco general para cada objetivo dentro del cual escoger nuestro instrumento más adecuado. 

\textcolor{red}{\textbf{PONER EJEMPLOS DE LO ANTERIOR:} queda más claro}

\clearpage
\section{Regulación: ¿Cómo?}
\puntosclave{
    La regulación mediante incentivos surge como respuesta a los problemas de la teoría tradicional para explicar fenómenos observados en la práctica regulatoria. 

    Entre sus principales rasgos se encuentra la introducción de problemas de información asimétrica y consideraciones normativas que pasaban por alto los modelos tradicionales (en particular la importancia de los incentivos de empresas y reguladores y su posible desviación de la maximización de los beneficios y bienestar social, respectivamente).
    
	El modelo pionero de la nueva teoría de la regulación económica es el de LAFFONT y TIROLE (1986), que introduce la existencia de una ventaja informativa por parte de la empresa regulada: el regulador no puede observar con exactitud su nivel de eficiencia (selección adversa) ni su esfuerzo (riesgo moral).
    
	En este contexto, lo óptimo para el regulador es ofrecer un menú regulatorio: un contrato con un precio relativamente fijo (fixed-prize) dirigido a las empresas más eficientes y otro que cubra los costes para las empresas menos eficientes (cost of service).  El regulador se enfrenta a un trade-off entre reducir la renta transferida a la empresa eficiente y reducir la ineficiencia introducida en el nivel de esfuerzo de la empresa ineficiente.
    
	El modelo puede extenderse, introduciendo fenómenos adicionales como auditorías de costes, problemas de inconsistencia dinámica por parte del regulador (ratchet effects), regulatory lag o la captura del regulador, entre otros.
    
	Estas nuevas teorías impulsaron la adopción de nuevos esquemas regulatorios como los precios máximos (price cap) o la compartición de beneficios (sliding scale).
}

Ahora que conocemos el cuándo y por qué hay situaciones en las que el policy-maker llave regula un mercado concreto, veámos cómo puede formularse de manera óptima. Aunque existen dos contextos diferentes en los que aplicaremos la regulación, nos centraremos en el caso de monopolios deseables.

En este contexto, ¿cuál es nuestro objetivo como reguladores? Para contestar a esta pregunta debemos únicamente recurrir a la justificación original de la intervención: el \textbf{precio}. Lo que nos interesará será buscar la eficiencia asinativa mediante la disminución del precio que otrosí tendría el producto en ausencia de regulación teniendo en cuenta la estructura productiva del oferente, que puede ser tal que la fijación del precio al coste marginal conduzca a pérdidas que provoquen el cierre del mercado. Por tanto, ¿qué debemos hacer?

\subsection{Objetivos: provisión e innovación}
Como todo en Economía, se trata de buscar un punto medio. Como reguladores queremos que el monopolista preste el servicio, puesto que de no hacerlo la pérdida social sería mayor que la apropiación del excedente que tendría lugar en un escenario de monopolio no regulado. No obstante, este servicio debe prestarse de forma tal que en la medida de lo posible se maximice el excedente neto del consumidor pues sino habría una pérdida social demasiado grande. 

Con el fin de resumir todo el abanico de objetivos que persigue un regulador, presentaremos los dos más importantes

\subsubsection{Precios RAMSEY-BOITEAUX}
Lo primero que persigue un regulador es la provisión del servicio al mayor nivel de eficiencia posible. Teniendo en cuenta que el primer óptimo no es posible (competencia perfecta), el regulador persigue entonces el segundo óptimo. Pero, ¿cómo se determina?

La aportación teórica más importante es el mecanismo de precios RAMSEY-BOITEUX. La esencia teórico-conceptual de este mecanismo es proporcionar el límite teórico de eficiencia cuando hay que financiar un coste fijo con precios distorsionantes. Partiendo de la consideración de que si fijamos precios igual a coste marginal, obtenemos eficiencia asignativa, pero el operador incurre en pérdidas. ¿Cuál es la mejor asignación factible dado un coste fijo inevitable? Considerando que el precio debe ser lineal.

Entonces, suponiendo que el regulador está plenamente informado, podemos plantear el siguiente problema de maximización:

\eqblock{
\begin{aligned}
    \boxed{\begin{aligned}
        &\max W(\mathbf p, \mathbf q)=\sum_{i=1}^I\left(\int_{0}^{q_i(p_i)}P_i^d(x)\mathrm dx\right)-\sum_{i=1}^Ic_iq_i\\[0.5em]
        &\text{s.a.}\quad \underbrace{\sum_{i=1}^I(p_i-c_i)q_i(p_i)}_{\text{\scriptsize mark-ups acumulados}}\ge{F}
    \end{aligned}}\\[1em]
    &\text{Solución}\\[0.5m]
    &\over\begin{aligned}
        \\
        &\mathcal L=\sum_i\left(\int_0^{q_i(p_i)}P_i(x) \mathrm dx\right)-\sum_ic_iq_i+\lambda\left[\sum_i(p_i-c_i)q_i-F\right]\\[0.5em]
        &\text{CPO:}\quad -q_i+\lambda \left[q_i+(p_i-c_i)\frac{\mathrm d q_i}{\mathrm d p_i}\right]=0\\[0.5em]
        &\overset{\text{\scriptsize Dividimos por $q_i$}}{\Longrightarrow}\qquad {\frac{p_i-c_i}{p_i}=\frac{1}{\varepsilon_{p_i,x_i}}\cdot \frac{\lambda}{1+\lambda}}\quad \text{donde,}\;\;\varepsilon_{x_i,p_i}=\frac{\mathrm dq_i}{\mathrm dp_i}\frac{p_i}{q_i}\\[2em]
        &\text{Reescribiendo,}\qquad \underset{\text{\scriptsize $R$ es el \textbf{número de RAMSEY}}}{\boxed{\frac{p_i-c_i}{p_i}=\frac{R}{\varepsilon_{p_i,x_i}}}}\qquad \text{si,}\;\;
        \left\{\begin{aligned}
            &F \gg 0\;\Rightarrow\; \mu \gg 0\\
            &F\sim 0\;\Rightarrow\; p_i\sim c_i
        \end{aligned}\right.
    \end{aligned}
\end{aligned}
}{Desarrollo analítico: Precios RAMSEY-BOITEAUX. Reculación: ¿Cómo?}

Por tanto, los precios RAMSEY-BOITEUX constituyen un paso intermedio entre la fijación de precios unilateral del monopolista y la competencia perfecta: se trata de un puente entre ambos en el que se distribuye de manera adecuada la carga de los costes hundidos entre los diferentes precios del mercado, reduciendo la carga sobre aquellos que tienen una elasticidad menor y aumentándola sobre aquellos que tienen una elasticidad mayor. No obstante, y lo que nos resulta más interesante, esta misma lógica se puede aplicar para un bien, para dos bienes o para todos los bienes que conforman la economía. 

\begin{specialblock}{Ejemplo fijación de precios RAMSEY-BOITEAUX: Monopolista multiproducto}
    Una forma más sencilla de resolver este problema en un contexto de dos productos es la condición de RAMSEY. Según RAMSEY, para minimizar las pérdidas de peso muerto, se deben aumentar los precios a las demandas/ofertas rígidas y elásticas en la misma proporción, en relación con los precios que se cobrarían en la solución de primer óptimo (precio igual al costo marginal).

    El área del East Bay de la región de la Bahía de San Francisco incluye Berkeley, Oakland, Walnut Creek y numerosas otras ciudades, además de algunas zonas no incorporadas. En esta área se proporcionan dos formas de transporte público: el servicio de autobuses de la compañía Alameda-Contra Costa (AC) Transit y el servicio ferroviario del sistema Bay Area Rapid Transit (BART). Un organismo regional de transporte, la Metropolitan Transportation Commission (MTC), coordina el servicio entre las distintas agencias de transporte de la región de la Bahía de San Francisco y ejerce una supervisión considerable sobre las tarifas de cada una de ellas. Como ocurre con la mayoría de los proveedores de transporte público, AC Transit y BART son monopolios naturales en el sentido de que su costo marginal es inferior a su costo medio en el intervalo relevante de producción.  

    \imagenfit{RAMSEY_pricing_Ejemplo.png}{RAMSEY prices para AC y para BART}{\href{https://eml.berkeley.edu/~train/regulation/ch4.pdf}{\textit{Ramsey Prices.} Berkley}}

    En consecuencia, fijar el precio de cada servicio a su costo marginal haría que ambas agencias perdieran dinero. Dada la autoridad de la MTC, la posibilidad de aplicar precios de Ramsey es factible en esta situación. Para efectos de fijación de precios y de cubrir los costos, las dos empresas de transporte pueden considerarse como una sola. Los precios de Ramsey para los dos servicios son aquellos que proporcionan el mayor excedente para los viajeros y, al mismo tiempo, permiten que los ingresos combinados de las dos agencias cubran sus costos combinados. La MTC podría administrar cualquier subvención cruzada que se requiriera con los precios de Ramsey:\footnote{%
        Los ingresos individuales de cada proveedor no necesariamente cubrirán exactamente los costos de ese servicio. Si cada proveedor fijara sus precios por separado, al costo medio, los ingresos de cada servicio cubrirían sus propios costos. El valor de la tarificación de Ramsey en esta situación es que permite un mayor excedente, porque implica solo una restricción sobre los precios (a saber, que los ingresos combinados cubran los costos combinados) en lugar de dos restricciones (a saber, que los ingresos de cada uno de los dos proveedores cubran sus propios costos).
    }  [...]  
    \begin{la}
        \item Las combinaciones de precios en las que los ingresos combinados igualan los costos combinados se representan como la curva A; esta curva es la parte relevante de la isocuanta de beneficio cero; y,
        \item La curva B representa las combinaciones de precios en las que estos productos son los mismos para autobús y ferrocarril.
    \end{la}   
    
    La intersección de las curvas A y B es la combinación de precios de Ramsey, porque en ese punto se satisfacen ambas condiciones (operar sin pérdidas) y (maximizar la utilidad total). Los precios de Ramsey son 2,42 centavos por milla para el autobús y 1,28 centavos por milla para el tren. A modo de comparación, el costo medio del viaje en autobús es de 2,0 centavos por pasajero-milla, y el costo medio de operación del tren es de 1,78 centavos. Debido a que el precio de Ramsey para el autobús excede el costo medio del servicio de autobús, mientras que el precio del tren es inferior al costo medio de operación de BART, el servicio de autobús en este caso estaría subvencionando al tren.\footnote{%
        En muchas aplicaciones de los precios de Ramsey, algunos costos se comparten conjuntamente en la producción de los bienes. Por ejemplo, se utiliza la misma capacidad de generación para producir electricidad para clientes residenciales y comerciales, aunque a los dos grupos se les cobre precios diferentes. En estos casos, no es posible calcular por separado el costo medio de cada servicio. Sin embargo, en el caso de AC Transit y BART, no se comparten costos, de modo que puede calcularse el costo medio de cada servicio.
    } 
\end{specialblock}

\paragraph{Implicaciones de Política Económica}
Desde una perspectiva de Política Económica, este mecanismo nos proporciona el precio óptimo teniendo en cuenta la necesaria financiación de un coste fijo que hace imposible el primer óptimo, obteniéndose un \textbf{second-best}. Este precio depende directamente del \textbf{número de RAMSEY} (entre 0 y 1) y sus implicaciones son profundas:
\begin{la}
    \item Primero, justifica una discriminación de precios eficiente desde el punto de vista del bienestar social, aun cuando los precios se alejen del costo marginal;
    \item Segundo, proporciona una base normativa para la regulación de monopolios naturales, mostrando que imponer precios iguales al costo marginal puede ser inviable financieramente; y,
    \item Tercero, introduce un criterio claro de equidad-eficiencia: la carga del financiamiento de los costos fijos recae más en quienes reaccionan menos a los precios, reduciendo así la distorsión total del sistema.
\end{la}   

En conjunto, conecta de manera rigurosa la teoría del bienestar, la elasticidad de la demanda y la regulación óptima de precios que dependerá de:
\begin{lnum}
    \item \textbf{Elasticidad de la demanda}. Cuánto más rígida sea la demanda del bien considerado mayor será el \textit{mark-up} sobre el precio pues mayor es el poder de mercado que ostenta el monopolista;  y,
    \item \textbf{Número de RAMSEY}. Cuánto mayor sea el valor de $\lambda$, que refleja el coste sombra de la restricción impuesto por los costes hundidos que enfrenta el monopolista: (i) a medida $\lambda\to 0$, el precio de RAMSEY-BOITEAUX se acerca al coste marginal (óptimo de primer orden); y, (ii) a medida que $\lambda \to\infty$, el precio se acercará al de monopolio, intuitivamente esto implica que la restricción es tan fuerte que la empresa regulada necesita el precio de monopolio para no incurrir en pérdidas. 
\end{lnum} 

Pese a que los precios de RAMSEY-BOITEUX proporcionan unos principios útiles para los reguladores en términos de minimización de la ineficiencia, su implementación práctica presenta varios problemas: (i) distributivos, que pueden surgir si no existen mecanismos para compensar a los perdedores con la aplicación de los precios de Ramsey; y, (ii) ausencia de otros fallos de mercado.

\subsubsection{Información asimétrica: esquemas de incentivos}
Ahora bien, los reguladores disponen de dos tipos de fuentes de información: (i) contable; y, (ii) demanda de mercado. Esto creo un problema evidente a la hora de implementar cualquier precio óptimo: la información asimétrica sobre los costes y el esfuerzo del operador. Es más, diferencia de la demanda que, en determinados casos, puede inferirse del comportamiento del mercado, los parámetros tecnológicos relevantes para la eficiencia productiva son privados del operador y no observables por el regulador. En este contexto, incluso si el regulador conoce el esquema de precios de segundo óptimo, su implementación efectiva es imposible sin un mecanismo que discipline el comportamiento de la empresa, ya que el operador tiene incentivos a inflar costes, reducir esfuerzo o explotar su ventaja informativa para apropiarse de rentas: no basta con fijar precios eficientes, es imprescindible diseñar contratos que alineen los intereses del operador.

Desde una perspectiva analítica, distinguimos entre dos tipos de restricciones informativas: el riesgo moral y la selección adversa. El riesgo moral se refiere a variables endógenas que no son observadas por el regulador. Utilizamos la etiqueta de \textbf{esfuerzo} de forma genérica para tales medidas discrecionales, que representa el número de horas de oficina que dedican los directivos de una empresa o la intensidad de su trabajo, pero que debe interpretarse de forma más amplia.\footnote{%
    Por ejmplo, algunos ejemplos de esfuerzo negativo (TIROLE y LAFFONT, 1989) son: (i) la asignación de ventajas por parte de los directivos; (ii) la indulgencia en actividades que privilegian su potencial profesional por encima de la eficiencia; (iii) el retraso de acciones desagradables; (iv) la compra de materiales y equipos a precios elevados; y/o, (v) el acaparamiento de ingenieros o máquinas que no son necesarios para los contratos actuales, pero que son útiles para obtener beneficios comerciales o para ganar contratos futuros son ejemplos de esfuerzo negativo.
} Por otro lado, la \textbf{selección adversa} surge cuando la empresa tiene más información sobre algunas variables exógenas. Es probable que esta asimetría de información incida sobre en las posibilidades tecnológicas de la empresa o en la dificultad para llevar a cabo determinadas tareas productivas.\footnote{%
    Por ejemplo, los contratistas de defensa suelen estar mejor informados que el Departamento de Defensa sobre el coste probable del desarrollo de un nuevo sistema de armas. Del mismo modo, varios autores han observado que AT\&T, y no la Comisión Federal de Comunicaciones, tiene la experiencia necesaria para pronosticar los costes de los servicios de telecomunicaciones, o que las empresas eléctricas saben más que las comisiones de servicios públicos.
}

La presencia de riesgo moral y selección adversa, y la consiguiente pérdida de control por parte del regulador, crean una demanda de recopilación de información.\footnote{%
    En la mayoría de los países, las empresas públicas son supervisadas periódicamente por organismos públicos. Del mismo modo, las empresas privadas reguladas están sujetas a controles por parte de organismos, y en ocasiones se les rechazan sus costes.
} Esta se proyecta en algunos casos como mecanismos de supervisión o revisón de las condiciones del contrato (riesgo moral) o mediante la necesidad de inducir la revalación por parte de la empresa que licita dicho contrato u opera dentro del mercado (selección adversa).

En definitiva, la regulación óptima pasa entonces de ser un problema puramente de asignación a uno de revelación y motivación. Este enfoque, desarrollado formalmente por LAFFONT y TIROLE, establece el segundo gran pilar de la teoría moderna de la regulación: \textbf{la eficiencia regulatoria no depende solo de las condiciones pactadas, sino de alinear los incentivos de las partes}.

\paragraph{Implicaciones de Política Económica}
Por tanto, en términos de Política Económica, lo esencial será determinar los esquemas de incentivos óptimos. Y, estos a su vez se materializarán a través de los dos instrumentos principales de los que dispone el regulados (que constituyen las dos vertientes de la regulación):
\begin{la}
    \item \textbf{Transferencias y/o precios de mercado}. Si el gobierno puede subvencionar (o gravar) a las empresas reguladas, es decir, si estas pueden recibir fondos públicos, el instrumento que debe vertebrar los esquemas de incentivos será la \textbf{transferencia};\footnote{%
        Las transferencias deben entenderse en un sentido amplío pues pueden adoptar varias formas: subvenciones directas, préstamos gubernamentales a tipos de interés bajos o sin obligación de devolución, garantías gubernamentales gratuitas cuando la empresa solicita préstamos en el mercado privado, transferencias de insumos públicos a precios deflactados, etc. Las transferencias también son bastante comunes cuando la empresa regulada es una empresa estatal.
    
    Por ejemplo, los ferrocarriles estadounidenses recibieron subvenciones tras su nacionalización en 1976. Del mismo modo, el Servicio Postal de los Estados Unidos ha recibido asignaciones. En el caso de empresas públicas como Electricité de France, que se supone que no reciben transferencias del gobierno pero que pueden obtener préstamos previa aprobación del Ministerio de Finanzas, se debe seguir razonando como si la empresa recibiera transferencias: cuando la empresa toma prestado 1 dólar, el resto del estado reduciría su propio endeudamiento en 1 dólar y aumentaría los impuestos en 1 dólar. En el margen de una política financiera óptima, el bienestar no se ve afectado por la reducción del endeudamiento y el aumento de los impuestos, y todo es como si la empresa recibiera 1 dólar del gobierno recaudado a través de los impuestos. Por último, las transferencias pueden tener lugar en algunos sectores en los que las empresas no son públicas. Un ejemplo de ello es el reembolso de Medicare.
    
    De la misma forma, estas transferencias a la empresa se producen necesariamente en la contratación de proyectos públicos, como la adquisición de armas: dado que el gobierno es el único comprador, los costes no pueden cubrirse mediante cargos a los clientes privados.
    } y,
    \item \textbf{Poder de los esquemas de incentivos}. Si la legislación no permite a los reguladores gravar a las empresas ni concederles subvenciones, el instrumento que vertebra los esquemas de incentivos será la \textbf{carga que puedan trasladar a los consumidores del bien regulado}. Las industrias en las que el gobierno tiene prohibido transferir dinero a las empresas reguladas a menudo aplican tarifas de dos partes [\authorfont{Ver Tema 3.A.16}].\footnote{%
        JOSKOW y SCHMALENSEE (1986) han conjeturado que la prima fija que pagan los consumidores desempeñaría un papel similar al de la transferencia (ausente) del gobierno.
    }
\end{la}

\subsection{Modelo: TIROLE-LAFFONT}
El modelo canónico de regulación es el de TIROLE-LAFFONT. Éste parte de la existencia de un monopolio regulado que sufraga los costes operativos a través de transferencias directas financiadas mediante impuestos distorsionantes. El regulador busca el esquema de incentivos que maximice el bienestar social agregado y garantice la buena conducta de los gestores empresariales: ¿cómo podemos alcanzar los precios RAMSEY-BOITEAUX mediante incentivos sin mermar el espíritu innovador y de reducción de costes que debería guiar la acción empresarial? Es decir, ¿cómo alinear los intereses de la empresa regulada con los intereses del regulador? 

Existen múltiples versiones de este problema. Nosotros recurriremos al caso de una empresa uniproducto de las características reseñadas. 

\subsubsection{Desarrollo formal: Empresa pública}
Ante una situación de monopolio público, el regulador toma el papel de Principal y los gerentes de esta empresa toman el papel de Agente. Para simplificar, consideraremos que existen dos niveles de eficiencia de la empresa: $\underline{\beta}\;,\overline\beta$. El regulador no sabe que nivel de eficiencia tiene esta empresa pero sí su probabilidad: $\mathrm P[\beta=\underline{\beta}]=\nu\;,\;\mathrm P[\beta=\overline{\beta}]=(1-\nu)$; por lo que se enfrenta a un problema de selección adversa. El regulador puede observar los costes medios de la empresa: $c=\beta-e$

Imaginemos que éstos se reintegran a la empresa atendiendo a un contrato adecuadamentemente diseñado (el que vamos a intentar encontrar). Además, los ingresos de la empresa son transferidos al regulador y  este utiliza impuestos distorsionantes (sobre el producto) para reintegrar los costes/esfuerzo, $t$. Con todo lo anterior, podemos formular el siguiente problema de maximización del regulador:

\eqblock{
\begin{aligned}
    &\boxed{\begin{aligned}
        &\max_{\{(\underline{t},\underline{q},\underline{C},\underline{c},\underline{U}),(\overline{t},\overline{q},\overline{C},\overline{c},\overline{U})\}}W=\nu\cdot W(\underline{t},\underline{q},\underline{C},\underline{c},\underline{U})+(1+\nu)\cdot (\overline{t},\overline{q},\overline{C},\overline{c},\overline{U})\\[1em]
        &\text{s.a.}\quad
        \left\{\begin{aligned}
            &\text{R.P:}\quad \overline{U}\ge0\;\;\;\text{y,}\;\;\;
            \underline{U}\ge 0
            \\[1em]
            &\text{R.C.I:}\quad 
            \left\{\begin{aligned}
                &\underline{U}= \underline{t}-\psi(\underline{\beta}-\underline{c})&&\ge \overline{t}-\psi(\underline{\beta}-\overline{c})&&(\text{\scriptsize eficiente no imite ineficiente})\\[0.5em]
                &\overline{U}= \overline{t}-\psi(\overline{\beta}-\overline{c})&&\ge \underline{t}-\psi(\overline{\beta}-\underline{c})&&(\text{\scriptsize ineficiente no imite eficiente})
            \end{aligned}\right.
        \end{aligned}\right.
    \end{aligned}}\\[0.5em]
    &\text{Donde,}\quad \underset{\text{\scriptsize Excedente total}}{W(U,q,e,\beta)}=\underbrace{\underset{\substack{\text{\scriptsize Excedente}\\\text{\scriptsize Bruto del}\\\text{\scriptsize Consumidor}}}{S(q)}+\underbrace{\lambda qP_x^d(q)}_{\substack{\text{\scriptsize Recaudación}\\\text{\scriptsize fiscal distorsionante}}}}_{\text{\scriptsize Lado de la \textbf{DEMANDA}}}-\underbrace{\underbrace{(1+\lambda)[\underbrace{(\beta-e)q+\psi(e)}_{\substack{\text{\scriptsize Coste variable de producción y}\\\text{\scriptsize coste del esfuerzo}}}]}_{\substack{\text{\scriptsize Costes variables de producción teniendo en cuenta}\\\text{\scriptsize financiación fiscal}}}-\underbrace{\lambda U}_{\substack{\text{\scriptsize Rentas}\\\text{\scriptsize informacionales}}}}_{\text{\scriptsize Lado de la \textbf{OFERTA}}}\\[1em]
    &\text{Solución}\\
    &\over
    \begin{aligned}
        \\&\text{Antes de $\mathcal L$ aglutinamos restricciones:}\qquad \left(
            \left.\begin{aligned}
                &\underline{U}= \underline{t}-\psi(\underline{\beta}-\underline{c})=\overline{t}-\psi(\underline{\beta}-\overline{c})=\\[0.5em]
                &=\overline{U}+\psi(\underline{\beta}-\overline{c})=\Phi(\overline{\beta}-\overline{c})=\Phi(\overline{e})
            \end{aligned}\;\;\right|\;\;\;\Phi(e)\equiv\psi(e)-\psi(e-\Delta\beta\right)\\[1em]
        &\text{\scriptsize $\Longrightarrow W=\nu[S(\underline{q})-(1+\lambda)[(\underline{\beta}-\underline{e})\underline{q}+\psi(\underline{e})]-\lambda\Phi(\overline{e})]+(1-\nu)[S(\overline{q})-(1+\lambda)[(\overline{\beta}-\overline{e})\overline{q}+\psi(\overline{e})]]$}\\[1em]
        &\underset{\text{\scriptsize Resolviendo}}{\text{CPO's}}:\;\;
        \left\{\begin{aligned}
            &\text{Eficiente:}\quad 
            \left.\begin{aligned}
                &\underline{q}=q^*(\underline{\beta}-\underline{e})\\[0.5em]
                &\psi'(\underline{e})=\underline{q}
            \end{aligned}\right\}
            &&\left(\substack{\text{\scriptsize Como las rentas de información dependen únicamente}\\\text{\scriptsize de $\overline{e}$, la cantidad y el esfuerzo de la empresa ineficiente}\\\text{\scriptsize es igual que en el escenario de información completa}}\right)\\[0.5em]
            &\text{Ineficiente:}\quad 
            \begin{aligned}
                &\overline{q}=q^*(\overline{\beta}-\overline{e})\\[0.5em]
                &\psi'(\overline{e})=\overline{q}-\left[\frac{\lambda}{1+\lambda}\right]\cdot \left[\frac{\nu}{1-\nu}\right]\cdot \Phi'(\overline{e})
            \end{aligned}
        \end{aligned}\right.
    \end{aligned}
\end{aligned}
}{Fijación de precios en monopolio uniproducto: El caso de dos tipos. Regulación en informaión asimétrica: Modelo de TIROLE-LAFFONT}

Por tanto, como vemos el problema de maximización es igual al de información completa (es decir, sujeto a restricción de producción/tecnológicaes), excepto por el coste esperado de las rentas informacionales $\nu\lambda\Phi(\overline{e})$. Como este nuevo término depende únicamente del esfuerzo empleado por los gerentes de la empresa menos eficiente ($\overline{\beta}$), los niveles de producción y esfuerzo de la empresa más eficiente ($\overline{\beta}$) son iguales que en información completa. Las CPO's relativas a óptimos de producción nos muestran que los precios óptimos para cada tipo de empresa es:

\eqblock{
\begin{aligned}
    &\text{Precios RAMSEY:}&&\boxed{\underline{p}=\overline{p}=p\Longrightarrow\quad \underset{\text{\scriptsize dado por la fórmula de RAMSEY}}{\frac{p-c}{p}=\frac{R}{\varepsilon_{x,p}}}}\\[0.5em]
    &\text{Rentas informativas y distorsiones:}&&\overline{U}\gg 0\quad \text{y,}\qquad \left[(\overline{e},\overline{q}, \overline{c})\leq\underbrace{(\overline{e}^*,\overline{q}^*,\overline{c}^*)}_{\text{\scriptsize Información completa}};\quad (\underline{e},\underline{q}\underline{c})=(\underline{e}^*,\underline{q}^*, \underline{c}^*)\right]
\end{aligned}
}{Precios RAMSEY. Regulación en informaión asimétrica: Modelo de TIROLE-LAFFONT}

Los precios no se ven distorsionados por la extracción de rentas o con el fin de generar incentivos: el sistema de incentivos viene completamente definido por la opción óptima del menú de contratos que el Agente atendiendo a la información privada que posee. Además, mientras el output y el esfuerzo de la empresa eficiente son iguales que en un escenario de información completa, disfruta de una renta informacional positiva por la información privilegiada: distorsión de incentivos. Por el contrario, se puede demostrar que el esfuerzo dedicado dado el tipo ineficiente es menor en condiciones de asimetrías informativas que en el escenario de información completa.

En definitiva, con dos tipos de empresas/agentes y transferencias gubernamentales, el óptimo exhibe la propiedad dicotómica de que el precio condicional al coste marginal no se ve afectado por la asimetría informativa. La regla o esquema de compensación se ve afectado por la presencia de información asimétrica originando que el coste medio del tipo ineficiente sea superior al de información completa (y el output y esfuerzo menor).

\subsubsection{Desarrollo formal: Empresa privatizada}
Una objeción frecuente a la rama de modelos anteriores como modelos de regulación es que ignoran las realidades institucionales al suponer que el Estado puede subvencionar o gravar a las empresas reguladas. Si bien es cierto que como modelos de adquisición/empresa pública no cabría objetar nada, ¿qué sucede cuando los ingresos de la empresa por las ventas a los consumidores deben cubrir cumplir con una restricción presupuestaria endogeneizada? ¿en qué se diferencian las conclusiones económicas de las que se obtienen cuando las transferencias son legales?\footnote{%
    Las recomendaciones sobre los incentivos son bastante similares: el coste oculto de los fondos públicos en toda la economía se sustituye por la pérdida marginal de eficiencia asociada al aumento de la prima fija, que expulsa del mercado a los consumidores marginales. Por el contrario, las recomendaciones sobre la fijación de precios son bastante diferentes.
}

He aquí la gran aportación de la Nueva Teoría de la Organización Industrial: las principales cuestiones de incentivos son sólidas frente a las restricciones legales sobre las transferencias, mientras que las cuestiones de precios no lo son.

\subsection{Implicaciones de Política Económica: NTOI}
La Teoría tradicional de la regulación frenó en el punto del que brotó la nueva Teoría de la regulación: los incentivos. 

La teoría de la regulación debe reflejar el entorno normativo. Aunque no es necesario que se ajuste a la práctica, debe ser coherente con las estructuras de información, las limitaciones y los instrumentos viables de las empresas y los reguladores.

Existen tres tipos de limitaciones normativas: informativas, transaccionales y administrativas y políticas. En la práctica, estas limitaciones impiden al regulador aplicar la política que prefiere (sea cual sea).

Desde esta perspectiva, tres son las razones principales que motivan la aparición de la NTOI: (i) la información asimétrica; (ii) la factibilidad o deseabilidad de los compromisos adquiridos (por ejemplo, por problemas de inconsistencia dinámica); y, (iii) la posible captura del regulador.


\subsubsection{Compromisos: Inconsistencia dinámica}
Las empresas reguladas y sus organismos o ministerios suelen interactuar repetidamente. En el ámbito de la contratación pública, los procesos de producción en varias etapas y los pedidos relacionados generan interacciones repetidas entre proveedores y compradores. Lo ideal es que estas relaciones se rijan por contratos amplios a largo plazo: si las dos partes pueden comprometerse a un contrato a largo plazo al inicio de su relación, el compromiso óptimo para el contrato en cada periodo es el contrato estático óptimo. Es decir, \textbf{lo óptimo para el regulador es comprometerse a no explotar la información obtenida al observar el rendimiento de la empresa}. El compromiso es crucial para este resultado, ya que el regulador querría extraer por completo la renta de la empresa a partir del segundo período, una vez que la empresa haya revelado su eficiencia en el primer período. Sin embargo, en la práctica, las relaciones regulatorias y de contratación pública suelen regirse por una serie de contratos a corto plazo, por dos razones:
\begin{lnum}
    \item \textbf{Marco legal}. En primer lugar, la mayoría de los países tienen prohibiciones legales sobre los compromisos a largo plazo; es decir, las administraciones actuales solo pueden vincular a las futuras en una medida limitada; y,
    \item \textbf{Contratación incompleta y/o inconsistencia dinámica}. La segunda motivación para la no vinculación es que la producción o las tecnologías futuras no pueden describirse perfectamente en un contrato hoy en día, y que los cambios en el entorno o el diseño pueden hacer que los contratos a largo plazo actuales pierdan su sentido. Esta es la motivación de la contratación incompleta.
\end{lnum}  

La falta de compromiso y la consiguiente negociación repetida entre las empresas y los reguladores tiene dos consecuencias normativas perversos. 

\paragraph{Enfoque normativo: \textit{Ratchett effects}}
En primer lugar, el regulador podría no compensar adecuadamente a la empresa por sus inversiones y, por lo tanto, esta se mostraría reacia a invertir. \textbf{Podemos seguir con algo de explicación aquí: TIROLE LAFFONT Cap.2}

En segundo lugar, como el regulador solo puede comprometerse con planes de incentivos de un período, en el segundo período el regulador diseñará su contrato óptimo en función de la información obtenida en el primer período contractual. Esto genera una distorsión muy potente en la conducta de la empresa: si produce a bajo coste en el primer periodo, el regulador podría deducir que los bajos costes no son tan difíciles de alcanzar y mañana ofrecer un contrato más exigente, apropiándose de las rentas informacionales. Es decir, la empresa pone en peligro las rentas futuras al ser eficiente.

Esto es lo que se conoce como \textit{\textbf{ratchett-effect}} o el efecto ratchet: el hecho de que la empresa esté preocupada por la expropiación de su renta informativa en el segundo período hace que la separación de tipos sea costosa y requiere la compensación adecuada en el primer periodo.

\paragraph{Enfoque positivo: \textit{Pooling equilibrium} y \textit{hit-and-run}}
Lo que nos conduce a un segundo enfoque valorativo: las implicaciones positivas de la interacción dinámica en el equilibrio. La NTOI demuestra que con un continuo de tipos de eficiencias potenciales y para cualquier sistema de incentivos del primer período, no existe un equilibrio de separación óptimo. Pero, ¿por qué? 

Este es uno de los resultados menos intuitivos. La causa principal de éste hecho y de la complejidad de los equilibrios que obtenemos se debe a que para inducir al tipo eficiente a revelar su tipo, el regulador debe concederle una gran renta en el primer periodo (por encima de su óptimo). Sin embargo, esta gran transferencia rompe el esquema de compatibilidad de incentivos, atrayendo al tipo ineficiente pues le resulta ventajoso seleccionar el contrato diseñado para el tipo eficiente (mimetización). Por tanto, a diferencia del marco estándar de un solo período presentado, donde solo la restricción de incentivos del tipo eficiente es vinculante, en un escenario dinámico donde sin compromiso, cualquier restricción de incentivos podría ser vinculante en el óptimo. 

De hecho, una vez mimetizado el tipo por parte de la empres ineficiente y obtenida la renta transfereida por el regulador, a la hora de diseñar el contrato del segundo periodo, el regulador, creyendo que la empresa es eficiente, deviene muy exigente, pero el tipo ineficiente abandona entonces la relación. Esta estrategia para el tipo ineficiente se conoce como \textit{\textbf{hit and run}} o \textit{coge el dinero y vete} y es la causa principal de la multiciplicidad y complejidad de los equilibrios en un escenario dinámiico.\footnote{%
    A pesar de este panorama sombrío de nuestro conocimiento sobre los sistemas de incentivos en una relación repetida regida por contratos a corto plazo, hay dos resultados sólidos en el caso de dos tipos que vale la pena destacar. En primer lugar, para factores de descuento pequeños, el contrato óptimo del primer período separa los dos tipos. Este resultado es intuitivo, pero no trivial: mientras que la empresa descuenta sustancialmente el futuro y solo necesita pequeños incentivos para revelar su información, el regulador también atribuye en el período un un valor bajo a disponer de información en el período dos. En segundo lugar, para factores de descuento elevados, el resultado regulatorio del primer período no puede distinguirse empíricamente de un resultado de agrupación pura (sin información).
}

\paragraph{Recomendaciones: Compromisos y renegociación}
Pero, si nuestro objetivo es regular un mercado con la intención de que este se perpetúe en un marco intertemporal de la manera menos ineficiente posible, pero a la ver resulta imposible (o no deseable) adquirir compromisos que nos conduzcan a este óptimo, ¿de qué nos sirve todo lo expuesto?

Si consideramos que la capacidad de compromiso de las partes se sitúa entre el compromiso intertemporal total y la ausencia de compromiso intertemporal estaríamos ante un escenario de \textbf{compromiso y renegociación}. Considerando el compromiso en el sentido de que las dos partes firman un contrato a largo plazo que se aplica mientras una de las partes desee que se aplique se obtiene que la restricción de incentivos solo es vinculante al alza (no hay hit and run) pero, sin embargo, muestra una efecto trinquete similar a la del caso sin compromiso.\footnote{%
    Por tanto, el caso de renegociación se asemeja técnicamente al caso de compromiso en que las restricciones CI se comportan correctamente: en el caso de dos tipos, solo la restricción IC del tipo eficiente es vinculante. En cierto sentido, la principal diferencia entre estos dos casos es la posibilidad de de evitar la estrategia hit and run, lo que permite al regulador ofrecer más incentivos al tipo eficiente para que se separe, al tiempo que evita que el tipo ineficiente imite al eficiente en el primer periodo.
} Además: (i) el equilibrio es separador solo si el factor de descuento es pequeño;\footnote{%
    La probabilidad de equilibrio de que el tipo eficiente se agrupe con el tipo ineficiente aumenta con el factor de descuento y converge a uno cuando las partes se vuelven muy pacientes.
} y, (ii) para un continuo de tipos, separar completamente los tipos es factible pero nunca óptimo para el regulador. De manera simplificada, la siguiente tabla muestra una comparación de los resultados obtenidos en cada caso:

\imagenfit{Comparacion_TiposCompromisos.png}{Tabla de comparación entre compromiso intertemporal total, compromiso con renegociación y asuencia de compromiso intertemporal}{\textit{A Tehory of Incentives in Procurement and Regulation}. Capítulo 10. TIROLE y LAFFONT, 1983}

\subsubsection{Captura del regulador: Óptimo social}
A esta lista de restricciones administrativas hay que añadir las restricciones políticas que pueden derivarse de las leyes del Congreso o de las legislaturas locales en los Estados Unidos (es decir, los comités encargados de supervisar la agencia reguladora) o del poder ejecutivo en otros países. Los políticos pueden influir en las decisiones de las agencias mediante el control de sus asignaciones presupuestarias y mediante amenazas de cambiar las responsabilidades de la agencia y de destituir a sus altos funcionarios o humillarlos durante audiencias públicas. Estas restricciones administrativas y políticas han sido ignoradas en gran medida tanto por la teoría tradicional como por la nueva economía regulatoria. Nuestra filosofía, tal y como se desarrolla en la última parte de este libro, es que dichas restricciones no son exógenas, sino que están impulsadas por restricciones informativas y transaccionales. Tienen que ver con el hecho de que los propios reguladores son agentes de otras partes (políticos y, más fundamentalmente, sus electores). La teoría del bienestar de las agencias sostiene que estas se crean y se les otorgan poderes sustanciales debido a la necesidad de obtener información sobre las compensaciones normativas. Sin embargo, los reguladores pueden tener sus propios objetivos y deben recibir incentivos para implementar los objetivos de los principales políticos. 1 0 Por lo tanto, la intención detrás de no permitir que una administración firme contratos a largo plazo con los proveedores parece ser limitar las externalidades entre administraciones: al no permitir que la administración actual vincule a las futuras, la ley reduce la eficiencia de los contratos, pero aumenta la rendición de cuentas de cada administración. Del mismo modo, la prohibición de las transferencias del gobierno a las empresas reguladas podría atribuirse al temor de que los reguladores abusaran de este instrumento. Los requisitos de procedimiento tienen por objeto limitar los acuerdos secretos entre la agencia y la industria u otros grupos de interés, generar información para los responsables políticos y permitirles reaccionar ante las políticas propuestas; son instrumentos para el control de las agencias. Por lo tanto, el estudio de la «regulación jerárquica», en lo que se refiere a la gobernanza de las estructuras reguladoras y políticas de varios niveles, debería ser un elemento central de la agenda de la nueva economía reguladora.

Por ello, es importante crear instituciones creíbles, con un marco procesal que limite las actuaciones del regulador y que los tribunales puedan dirimir los potenciales conflictos con las empresas. Si se amplía el horizonte temporal, el regulador puede tener interés en cumplir sus compromisos para mantener una buena reputación. Un elemento clave para que la reputación pueda sostener el equilibrio es que el regulador esté desvinculado del ciclo político, de ahí la tendencia a crear reguladores independientes. No obstante, como señala JOSKOW (2007), en la realidad ninguna agencia reguladora puede ser del todo independiente del poder político (los miembros de sus consejos suelen ser nombramientos políticos y aunque no puedan ser despedidos sí que pueden no ser renovados, los reguladores se encuentran sometidos al control parlamentario, pueden surgir problemas de falta de fondos y de personal, etc.)

\subsubsection{Información asimétrica: Esquemas de incentivos}
Una vez mencionadas estas dos grandes aportaciones de la NTOI, volvamos al núcleo teórico de esta nueva vertiente: la información asimétrica. 

Como los reguladores no pueden basarse en contratos regulatorios que dependen de información que solo posee la empresa, las restricciones informativas limitan la eficacia del control de las industrias por parte de éstos y los esquemas de incentivos deben diseñarse de manera tal que garanticen que la empresa esté dispuesta a participar en el proceso de producción. 

Ahora bien, aunque hemos hablado con frecuencia sobre el esquema de incentivos, ¿cuál es? ¿cómo se materializa? La razón de su omisión es muy importante y se debe a que el esquema de incentivos óptimo dependerá de las características del sector/mercado regulado. Esto es clave: la NTOI no discrimina planes, los encuadra dentro del escenario en el que constituyen la forma óptima de regulación. 

\eqblock{
\begin{aligned}
    &\underset{\text{\scriptsize Esquea de incentivos lineal}}{\text{Linear scaling plan:}}\quad {\boxed{t=a+(1-b)C}}\qquad \text{donde,}\;
    \left\{\begin{aligned}
        &t&&\equiv &&\text{Renta obtenida por la empresa}\\
        &(1-b)&&\equiv &&\underset{\text{\scriptsize Determina el \textbf{poder} del esquema de incentivos}}{\text{Renta obtenida por reembolso}}\\
        &a&&\equiv&&\text{Renta por cuota fija definida}
    \end{aligned}\right.\\[1em]
    &\text{Teniendo en cuenta que,}\;\;
    \underset{\substack{\\[0.5em]\\\text{\scriptsize Cualquier contrato con una pendiente $b\in(0,1)$ se definne como \textbf{contrato de incentivos}}}}{\left\{\begin{aligned}
        &b\to 1\Longrightarrow\text{\textit{High-powered} incentive schemes}\\[0.5em]
        &b\to 0\Longrightarrow\text{\textit{Low-powered} incentive schemes}\\[0.5em]
    \end{aligned}\right.}\\[2em]
    &\text{\textbf{C.-} El rendimiento ($-C$) de la empresa esta correlacionado positivamente con $b$, la pendiente, y con el pago fijo $a$}.\\[1em]
    &\Longrightarrow\qquad \underset{\substack{\text{\scriptsize No interpretar como relación causal: contratos high-powered inducen eficiencia (X);}\\\text{\scriptsize sino como análisis descriptivo: empresas eficientes prefieren high-powered ($\checkmark$)}}}{\boxed{\frac{\mathrm dC}{\mathrm d\beta}>0;\;\;\frac{\mathrm db}{\mathrm d\beta}<0;\;\;\frac{\mathrm da}{\mathrm d\beta}<0}}
\end{aligned}
}{Diseño óptimo de los planes de incentivos: Esquemas lineales. Regulación en presencia de información aismétrica: Implicaciones de Política Económica}

La regulación óptima consiste en que el regulador debe ofrecer un \textbf{menú de contratos lineal en costes}.\footnote{
Los contratos del mundo real suelen ser lineales, pero algunos tienen características no lineales, como un límite máximo en las transferencias del gobierno o una garantía de que la empresa no perderá dinero. En este escenario la linealidad se da por tramos. Otro ejemplo de contrato lineal por tramos es una opción sobre acciones para directivos, que recompensa al directivo de forma lineal en función del valor de la empresa por encima de un umbral determinado.

De todas formas, la linealidad no se mantiene en caso de aversión al riesgo por parte de la empresa regulada o en situaciones de regulación dinámica. Cuando el esquema no es lineal, el poder del esquema de incentivos se identifica con el nivel de esfuerzo que esta emplea en el equilibrio.
} Este menú de contratos debe estar diseñado para diferentes tipo de empresa, con el objetivo de discriminar o inducir la auto-selección de diferentes tipos potenciales de empresa de la misma manera que un monopolista discrimina en precios entre consumidores con diferentes valoraciones respecto a cantidad o calidad.\footnote{%
    Además, independientemente de si la empresa enfrenta costes intrínsecamente altos o intrínsicamente bajos debe disfrutar de una renta no negativa. 

En el peor escenario, la empresa podría reducir su esfuerzo/inversión en reducción de costes por debajo del óptimo social y producir a un coste más elevado del que sería incluso en caso de que la empresa fuese muy ineficiente. Este recorte proporciona más utilidad a la empresa de la que hubiese disfrutado siendo ineficiente, y por ello una renta estríctamente positiva. Aunque la esencia de la renta es la posibilidad de reducir esfuerzo, vemos como en el escenario de regulación óptima esta renta debe ser garantizada; KHAN destaca que la información asimétrica da origen a rentas indeseadas que remuneran la suerte y desarrollos externos favorables.
} Gráficamente:

\textcolor{red}{\textbf{AQUI PONER EL SLOPE DEL ESQUEMA DE INCENTIVOS Y LO DE LA INTRODUCCIÓN MARCADO CON EL PAPEL BLANCO}}

ACABAR CON LA TABLA DE LOS DIFERENTES "TIPOS DE CONTRATO EN EL MUNDO REAL" PARA TENER UNA LINEA ADECUADA PARA EL SIGUIENTE APARTADO!!! ES MUY BUENA IDEA

\textbf{Pág. 40}

Por último, cabe incidir en la conclusión final de la NTOI relativa a la motivación de los diferentes tipos de contratos. Citando a TIROLE-LAFFONT (A Theory of Incentives in Procurement and Regulation, 1989):

\textit{Es importante no interpretar la proposición como una relación causal. De lo contrario, se podría concluir que los planes de incentivos de alto rendimiento ($b$ alto) son mejores porque inducen un mejor rendimiento. Si bien la segunda afirmación es correcta ((...) inducen un mejor rendimiento), la primera ignora la conveniencia de la extracción de rentas. La selección óptima de la tecnología de la empresa da como resultado, en la muestra, buenos rendimientos y planes de alta potencia, junto con malos rendimientos y planes de baja potencia. [...] La base de la conclusión es que los tipos más eficientes eligen planes de alta potencia porque producen a un coste menor y, por lo tanto, no les importa asumir su coste a cambio de una cuota fija elevada.}


\clearpage
\section{Regulación: Aplicaciones prácticas}
\puntosclave{
    La existencia de varias empresas en un mercado regulado puede servir para reducir la desventaja informativa del regulador (competencia por comparación o yardstick).
    
	Otra alternativa a la regulación tradicional es la subasta de una concesión para explotar un servicio en régimen de monopolio. Este sistema permite introducir competencia ex ante entre las empresas participantes en la subasta.
    
	Un caso particularmente interesante es la regulación en mercados recientemente liberalizados. Muchos de ellos son industrias de red en las que existe un componente competitivo y otro con rasgos de monopolio natural.
    
	En estos casos de mercados verticalmente relacionados, hay que valorar si se debe permitir que el antiguo monopolista siga prestando servicios en el segmento competitivo y, en caso afirmativo, cuál debería ser el precio óptimo para los nuevos competidores que quieran acceder al segmento no competitivo (la red).
    
	Los mercados de varias caras (multi-sided markets) son uno de los temas más estudiados estos últimos años por la teoría de la regulación, debido al surgimiento de grandes plataformas digitales. A nivel de la UE, se ha aprobado la DMA para tratar de corregir las consecuencias adversas de estos mercados sobre el bienestar social.
}

Ahora que hemos visto el marco teórico de la NTOI, sus principales aportaciones, motivaciones e implicaciones en términos de Política Económica, podemos pasar a aterrizar estos resultados y ver como se proyectan en el mundo práctico de la regulación. 

Dos son los objetivos de este último apartado: (i) ver como la NTOI se proyecta en los esquemas regulatorios más comunes, presentando sus características, ventajas y limitaciones así como algunos ejemplos prácticos; y, (ii) levantar la restricción impuesta sobre la presencia ex ante de empresas de determinada característica y ver cómo la competencia del mercado puede reducir las asimetrías informativas.

\subsection{Esquemas de incentivos}
Tal y como hemos visto, los planes de incentivos pueden analizarse atendiendo a su poder, es decir, el vínculo entre la transferencia de la empresa del gobierno y/o los precios de la empresa y su rendimiento en términos de costes o beneficios. De forma resumida, los esquemas más utilizados en la actualidad son:

\imagenfit{Esquemas_Practicos.png}{Poder de los esquemas frecuentemente utilizados}{}

\subsubsection{\textit{Low-powered incentive schemes}}
Históricamente, el método más importante para regular mercados ha sido mediante esquemas de incentivos de bajo poder. Estos se caracterizan porque la empresa no es \textit{residual claimant}: los costes realizados son recuperables y el regulador, el Estado o los consumidores absorben el riesgo. El vínculo entre desempeño en costes y remuneración es débil, por lo que los incentivos a reducir costes y a innovar son mínimos, pero a cambio se garantiza viabilidad financiera, estabilidad del servicio y bajo riesgo de quiebra. 

\paragraph{\textit{Cost-plus}: Transferencias por costes}
El ejemplo paradigmático es el contrato cost-plus. Bajo este esquema, el regulador reembolsa a la empresa los costes efectivamente incurridos y añade un margen fijo o una tasa de retorno preestablecida. Este tipo de contratos ha sido ampliamente utilizado en defensa, obra pública y grandes proyectos de I+D, especialmente en Estados Unidos bajo el marco del Federal Acquisition Regulation (FAR). 

Su ventaja es clara: permite ejecutar proyectos con alta incertidumbre tecnológica o de costes sin exigir primas de riesgo elevadas. Sin embargo, la desventaja es estructural: la empresa tiene incentivos a inflar costes, sobredimensionar inversiones y relajar el esfuerzo, fenómeno conocido como cost padding. 

\paragraph{\textit{Cost Of Service}: Precios medios}
El análogo funcional del cost-plus sin transferencias es la regulación por coste del servicio (cost-of-service regulation o rate-of-return regulation). Este esquema fue dominante en la regulación de electricidad, gas y telecomunicaciones en Estados Unidos durante gran parte del siglo XX. El regulador fija tarifas que permiten a la empresa recuperar todos sus costes operativos y de capital, más una tasa de retorno razonable sobre la base de activos. Ejemplos clásicos incluyen la regulación de las electric utilities por las Public Utility Commissions estatales y la regulación de AT\&T antes de la liberalización del sector.

No obstante, un aspecto crucial para determinar las propiedades incentivadoras de la regulación cost of service es el período de tiempo durante el cual se fijan los precios de producción. Dado que los precios son rígidos y la empresa es un reclamante residual por sus ahorros de costes entre revisiones, la práctica de la regulación COS difiere de su espíritu de coste incrementado, que exige que los precios sigan los costes de forma continua. De hecho, con un desfase regulatorio infinito, la regulación COS tendría un carácter de precio fijo. En la práctica, las revisiones reglamentarias son endógenas y suelen ser iniciadas por la empresa regulada o por la agencia reguladora. Muy a menudo, una empresa que se enfrenta a un aumento de los costes de los insumos solicita a la agencia permiso para subir sus precios. 

Al celebrarse una audiencia, se recopila información sobre los costes y los ingresos durante un periodo de prueba de un año, por ejemplo. La presentación de tarifas, las audiencias reglamentarias y la toma de decisiones por parte de las agencias son largas. Los aumentos de precios cuentan con una fuerte oposición por parte de los clientes residenciales e industriales. Los precios son bastante rígidos tanto al alza como a la baja.


\subsubsection{Intermedios: Contratos de incentivos}
Tras el desencanto creciente de este tipo de contratos la atención se centró en formas alternativas,\footnote{%
    En el sector de las telecomunicaciones, la Oficina de Telecomunicaciones del Reino Unido adoptó límites máximos de precios para British Telecom en 1984, y en 1989 la Comisión Federal de Comunicaciones de los Estados Unidos hizo lo propio con AT\&T. En el sector energético, en 1986 se adoptaron límites máximos de precios para British Gas, mientras que en los Estados Unidos se implementaron diversas formas de regulación incentivadora para las empresas de servicios eléctricos. Otro ejemplo de experimentación con nuevos mecanismos de incentivos es el de la Comisión de Comercio Interestatal de los Estados Unidos, que permite a los ferrocarriles fijar los precios que deseen, con la restricción particular de que los ingresos totales no deben superar los costes totales en un horizonte de cuatro años.
} que proliferaron enormenete durante la década de los años 80's aunque el concepto era antiguo.\footnote{%
    Por ejemplo, los planes de escala móvil que ajustan los precios a la baja cuando la tasa de rendimiento de la empresa supera la tasa objetivo y en los que el ajuste es parcial para permitir que la empresa conserve parte de los beneficios obtenidos; se adoptaron en 1925 para la Potomac Electric Power Company (tras algunos ajustes, el plan se abandonó en 1955 debido a los efectos de la inflación). El precio se redujo con el fin de reducir los beneficios excesivos en un 50\% cada año sucesivo (por el contrario, una tasa de rendimiento inferior al objetivo daba lugar a un ajuste completo del precio para evitar pérdidas de ingresos). 

Otros ejemplos se han registrado en los informes del Laboratorio de Energía del MIT (1984) y la FERC (1989). En el caso de la New York Telephone Company, el plan de 1986 estipulaba que los requisitos de ingresos se ajustarían en una cantidad equivalente a $(r - 15\%) / 2$ si $r >15\%$ (ajuste a la baja), a $(13\% - r) / 2$ si $r < 13\%$ (ajuste al alza), y no se ajustaría si $13\%<r<15\%$, donde $r$ es la tasa de rendimiento obtenida. 

Muchos planes de incentivos para las empresas eléctricas vinculan los beneficios de las empresas a sus ahorros/sobrecostes de combustible, los costes de construcción, la tasa de calor de la planta y el factor de capacidad de la planta.
} Los contratos de incentivos constituyen un esquema de poder intermedio y puede adoptar multitud de formas diversas. 

Los esquemas de incentivos intermedios surgen históricamente como respuesta al doble problema que plantean los esquemas extremos. Por un lado, los contratos de bajo poder (cost-plus, rate-of-return) generan escasos incentivos a la eficiencia y favorecen la inflación de costes; por otro, los esquemas de alto poder trasladan demasiado riesgo al operador y requieren condiciones institucionales y de información que no siempre están presentes. Los contratos intermedios buscan precisamente un equilibrio entre incentivos y aseguramiento, repartiendo de forma explícita los riesgos y los beneficios de las desviaciones de costes entre el regulador y la empresa. En términos conceptuales, el regulador no renuncia al control de costes, pero tampoco los reembolsa íntegramente: introduce una regla de sharing que suaviza el trade-off entre eficiencia productiva y rentas informacionales.

\paragraph{Compartir costes}
En los contextos donde las transferencias están permitidas, este enfoque se materializa en los denominados contratos de incentivos (incentive contracts). En ellos, el pago total al operador se compone de una transferencia base y de un componente variable que depende del desempeño en costes o beneficios. En estos esquemas, el regulador o comprador se compromete a cubrir una fracción predeterminada de los costes realizados, mientras que la empresa asume el resto. Cuanto menor es la fracción cubierta por el regulador, mayor es el poder del incentivo. La ventaja principal es que se mitiga el riesgo extremo del precio fijo cuando los costes son inciertos, sin caer en la neutralización total de incentivos del cost-plus. La contrapartida es clara: la empresa solo internaliza parcialmente el beneficio de reducir costes, por lo que el esfuerzo inducido es inferior al socialmente óptimo. Este tipo de contratos ha sido extensamente analizado en la literatura de contratos incompletos y procurement público.\footnote{
\href{https://www.nobelprize.org/uploads/2018/06/popular-economicsciences2016.pdf}{\textit{Contract Theory}. (NOBEL) THE PRIZE IN ECONOMIC SCIENCES 2016}
}

No obstante, cabe señalar que la regulación por incentivos puede aplicarse a otras medidas de rendimiento verificables además del coste. Al igual que la compensación en los contratos de defensa se basa en el cumplimiento de los plazos o en diversas dimensiones de calidad de los proyectos de adquisición, los ajustes de precios en los contratos regulados se refieren a las tasas de interrupción del servicio, los programas de conservación de energía para los consumidores, los servicios al consumidor, etc.

\paragraph{Compartir beneficios}
Un segundo ejemplo clave son los contratos de compartición de beneficios (profit sharing o sliding-scale regulation), históricamente relevantes en la regulación de utilities en Estados Unidos. Bajo estos esquemas, los beneficios de la empresa se comparan con un nivel de referencia, y las desviaciones —positivas o negativas— se reparten entre la empresa y los consumidores según una regla preestablecida. El objetivo es explícito: alinear los incentivos del operador con los del regulador, permitiendo que la empresa se beneficie de mejoras de eficiencia, pero limitando la acumulación de rentas excesivas. 

Estudios empíricos sobre regulación eléctrica y ferroviaria muestran que estos mecanismos mejoran el control de costes respecto al rate-of-return, aunque generan incentivos más débiles que los price caps puros y requieren una contabilidad regulatoria compleja.

La importancia de los contratos intermedios no reside en que sean teóricamente óptimos en un sentido abstracto, sino en que aproximan soluciones óptimas bajo fuertes restricciones institucionales: aversión política al riesgo, dificultades para medir calidad, oposición a quiebras de operadores esenciales y limitaciones legales a las transferencias. Por ello, siguen siendo predominantes en sectores donde la incertidumbre tecnológica es elevada o donde la continuidad del servicio es prioritaria, como transporte ferroviario, agua o determinadas concesiones energéticas. En suma, estos contratos ilustran un mensaje central de la NTOI: en regulación, la cuestión no es si introducir incentivos, sino cuánto poder darles y cómo repartir el riesgo, una decisión inevitablemente dependiente del contexto institucional y tecnológico.

\subsubsection{\textit{High-powered incentive schemes}}
Un esquema es high-powered cuando la empresa se aproxima a ser residual claimant: su beneficio aumenta casi uno-a-uno con las reducciones de coste (y, simétricamente, asume una parte relevante del riesgo). El regulador se ebenficia de eficiencia productiva sustituyendo el control detallado de costes por una regla simple (precio o pago) y acepta que la empresa retenga parte de las ganancias de eficiencia durante el periodo regulatorio/contractual. En términos prácticos, estos esquemas reducen el problema de \textit{cost padding} típico del cost-plus, pero introducen dos tensiones inevitables: (i) riesgo y prima de riesgo (si el operador asume shocks de coste/demanda), y (ii) riesgo de degradación de calidad si la calidad no está bien medida/penalizada.\footnote{%
    \href{https://economics.mit.edu/sites/default/files/2022-09/Incentive%20Regulation%20in%20Theory%20and%20Practice%20Electric%20Transmission%20and%20Distribution%20Networks%20%28revised%29.pdf?utm_source=chatgpt.com}{\textit{INCENTIVE REGULATION IN THEORY AND PRACTICE:
ELECTRICITY DISTRIBUTION AND TRANSMISSION NETWORKS}. JOSKOW, 2006
}
}

\paragraph{\textit{Fixed-priced}: Transferencias definidas}
Desde una perspectiva dinámica a menudo se afirma que los sistemas de incentivos high-powerd son poco prácticos en el ámbito de la contratación porque es probable que se renegocien. Sin embargo los planes de incentivos de este tipo sí que producen una producción eficiente y, en consecuencia, están menos sujetos a renegociación que otros planes diseñados para limitar los beneficios de la empresa. De hecho, si se compara el compromiso total intertemporal con el compromiso y la renegociación, la amenaza de renegociación induce a planes de incentivos más potentes en el segundo periodo abjo este esquema de incentivos. Por tanto, si la creencia convencional tiene algún mérito, debe provenir de factores que no se han tenido en cuenta en nuestro modelo, como la incertidumbre de los costes ex post, combinada con la quiebra o razones políticas para que el regulador no se comprometa ni siquiera con el contrato original.

Un instrumento claro de este tipo de esquemas en la contratación son los contratos de precio fijo y cerrado, donde el contratista asume plena responsabilidad por costes y beneficio/pérdida, lo que genera incentivo a controlar costes y reduce carga administrativa. Ejemplos típicos de su utlización son la contratación pública de bienes/servicios estandarizables (defensa, suministros, obras con alcance definido) bajo FAR/DoD y equivalentes europeos; y, dentro de empresa pública, contratos internos con presupuesto cerrado más métricas de desempeño.

\paragraph{\textit{Prices-cap}: Precios máximos}
Cuando el regulador no puede transferir (ni subsidiar) a la firma, el instrumento high-powered más usado es el price cap: se fija una senda de precios máximos por un periodo, permitiendo que la firma retenga ganancias si reduce costes por debajo de lo anticipado. Además, normalmente los precios máximos se ajustan de acuerdo con una fórmula pre-especificada con carácter anual, siendo la más conocida el $P_t=P_{t-1}+(IPC-X)$, por el cual el incremento anual de los precios no puede superar el incremento de un índice de precios exógeno y no manipulable por la empresa, menos un factor $X$ que refleja las ganancias potenciales de productividad. A largo plazo y con una periodicidad prestablecida  también pueden producirse modificaciones en los precios máximos, las cestas y los factores de ajuste.

En Reino Unido, el mecanismo RPI–X se implantó primero en telecomunicaciones con British Telecom tras la privatización (desde 1984, con controles tipo RPI–3\% y luego RPI–4.5\% en periodos tempranos), y se convirtió en referente internacional. En EE. UU., la FCC adoptó price caps para grandes carriers a comienzos de los años 90's como alternativa al rate-of-return. En España los precios máximos han sido sobre todo utilizados para regular las llamadas telefónicas (sector de las telecomunicaciones) y en la actualidad se aplican a las tasas aeroportuarias que AENA cobra a las aerolíneas.

Evidentemente, presenta una serie de ventajas por su propia naturaleza: (i) incentivos fuertes a eficiencia con menor necesidad de abrir los libros cada año; (ii) introduce un sustituto regulatorio de la presión competitiva (la empresa se beneficia de innovaciones/reducciones de coste durante el periodo); y, (iii) puede combinarse con incentivos de calidad y revisiones periódicas para limitar degradación de servicio. La evidencia y la discusión aplicada en redes (electricidad/gas/agua y telecoms) enfatizan precisamente estos trade-offs y la necesidad de complementar el cap con métricas de calidad. No obstante, está sujeto a una serie de limitaciones como: (i) problemas de calibración del precio y de $X$ (productividad esperada): si es demasiado exigente, riesgo financiero e inversión; si es laxo, rentas excesivas; (ii) incentivo a reclasificar costes o desplazar inversión/opex de forma estratégica si el modelo de revisión lo premia; (iii) riesgo de excluir consumidores o distorsionar estructura tarifaria si el cap se aplica de forma agregada sin cuidado (aquí es donde pricing se vuelve menos robusto que incentivos en la NTOI). \footnote{%
    Además, no relacionados con la determinación de los precios máximos elegidos ni el periodo, se encuentran:
\begin{lnum}
    \item Mayor grado de riesgo para la empresa regulada. Ya que: (i) por un lado, el riesgo regulatorio implica que si la empresa logra altos beneficios durante el periodo de vigencia, el regulador puede tener la presión política de intervenir y revisar los precios con el fin de confiscar los beneficios obtenidos y revertirlos a los consumidores vía menores precios. Sin embargo, con este tipo de intervención se reduce el incentivo de las empresas a buscar las máximas ganancias de eficiencia; y, (ii) por otro lado, si la incertidumbre respecto a los costes de la empresa es elevada y no se establecen mecanismos de ajuste adecuados, pueden producirse periodos con fuertes pérdidas que comprometan la viabilidad de la empresa. Esto puede llevar a que, en valor esperado, haya que fijar mayores precios que con una regulación basada en el coste del servicio, dando lugar a ineficiencia asignativa; y,
    \item Efectos sobre la calidad del servicio. La empresa regulada, viéndose incapacitada a subir los precios, puede optar por reducir la calidad, con lo que aumenta el precio por unidad de calidad. 
\end{lnum}

Efectos sobre la variedad de productos. Este mecanismo puede generar un incentivo a reducir la cartera de productos, especialmente aquellos que no cubran los costes incrementales o cuyo crecimiento de la productividad sea menor que $X$.
}

\subsection{Competencia: reduccion de asimetrías informativas}
Hasta ahora hemos supuesto que el regulador se ve restringido a actuar sobre una empresa o un conjunto definido de empresas. Pero este no es un supuesto que siempre su cumpla. Es evidente que, por ejemplo, en las industrias red el proveedor es único y titular de la infraestructura que determina su condición de monopolista. Por tanto no hay opción a considerar la posibilidad de que exista competencia por el mercado (de hecho, esto probablemente genere mayores ineficiencias).

Sin embargo, algunos mercados regulados están explícita o implícitamente abiertos a la competencia. Por un lado, existen sectores regulados no caracterizados como monopolio natural y sin que la competencia genera externalidades negativas que deba impedir el regulador. Por otro lado, la titularidad de la infraestructura o condición que determina la condición de monopolista puede ser de titularidad pública, pero el \textit{policy maker} puede optar (por voluntad u obligación) a delegar la provisión del bien/servicio a la gestión privada; situación en la cual el sector público podría escoger entre potenciales monopolistas a través de un proceso competitivo. 

En ambos casos es fácil ver la ventaja que emerge para el regulador: las asimetrías informativas se reducen debido a la competencia. De hecho, veremos brevemente como el regulador puede explotar esta ventaja realizable de la mejor manera posible.

\subsubsection{Competencia en el mercado: \textit{yardstick competition}}
Imaginémos que existen varias empresas operando en diferentes areas geográfica, dentro de un mismo mercado regulado, sea el de la generación eléctrica. El regulador puede reducir su desventaja informativa fijando el precio de cada empresa a partir de los costes medios del resto de empresas reguladas. 

Este mecanismo proporciona una solución al problema de revelación correcta de costes de las empresas reguladas, si son varias, en caso de que: (i) el regulador carezca de información específica sobre los mismos; y, (ii) no existan comportamientos colusorios.\footnote{%
    En general, la colusión se dificulta si: (i) el regulador puede fijar cotas superiores a los valores de los costes de cada empresa (derivadas de un conocimiento técnico del sector); y, (ii) el número de empresas del sector es elevado.
} 

Por tanto, la empresa no obtendrá beneficio alguno por proporcionar información errónea ya que si declara costes mayores de los verdaderos beneficiará a sus competidores, pero no a ella misma.  Como ventaja añadida, esta técnica regulatoria induce a la reducción eficiente de costes, ya que, si sus costes reales se sitúan por encima de la media de sus pares del sector, sólo se le retribuirá con arreglo a los costes del sector.

\subsubsection{Competencia por el acceso: concesiones}
Cuando existen varias empresas potenciacles para un mercado con rasgos de monopolio natural, conceder a una empresa la explotación exclusiva de un monopolio por un periodo de tiempo determinado mediante una subasta es el mecanismo óptimo para el regulador.\footnote{%
    Por ejemplo, la concesión se podría conceder a la empresa que ofrezca el servicio a un precio menor, aunque las autoridades podrán incorporar criterios adicionales. La idea original es de DEMSCTZ (1968). Su formulación más completa, que incluye una aplicación, se debe a POSNER. (1972). Para una crítica véase WILLIAMSON (1976); JOSKOW (2007), págs. 1267-1270 resume la literatura relacionada, con ejemplos prácticos en EE.UU.

En España, un sector en el que se aplica este modelo de intervención pública es el del transporte interurbano de viajeros en autobús. Para ver los detalles de la regulación y los problemas asociados a su aplicación práctica. Véase CNMC (2022).
} 

La idea subyacente es que se puede alcanzar un resultado similar al de los mercados impugnables (precio igual a coste medio) sin que exista competencia dentro del mercado, ya que la subasta produce competencia ex ante y la periodicidad de la misma competencia ex post. 

\paragraph{Eficiencia de la subasta}
Lo razonable es suponer que las empresas presentarán ofertas en función de sus costes medios más un margen razonable de beneficio (lo cual es seguro si existen rendimientos crecientes a niveles iniciales de producción), en cuyo caso nos encontraremos con precios subóptimos por determinarse en función del coste medio.

Podría pensarse en exigir ofertas con tarifas más complejas, pero en ese caso su valoración por parte de las autoridades sería más complicada. Piénsese, por ejemplo, en el tema de comparar distintas tarifas en dos partes que difieren en: 
\begin{la}
    \item La cuota de instalación (y, por lo tanto. en las clases de consumidores excluidas); y,
    \item El precio unitario (y, por lo tanto, en el grado de eficiencia).
\end{la}

En este caso, la dificultad de elección entre las distintas ofertas residiría en que cada oferta implicaría la eliminación de distintos grupos de consumidores del acceso al bien por la vía de la cuota fija. 

En general, el problema de la valoración de las ofertas radica en el hecho de que, si el sector público no presta el servicio o realiza la producción del bien, tendrá menos información que los oferentes particulares respecto a los costes y la demanda de mercado, por lo que su valoración de los términos en que se producen las ofertas será muy deficiente. Por otra parte. en el caso de prestación de servicios, resulta complejo comparar ofertas alternativas que afectan a un conjunto de variables tales como precios, calidades del servicio, frecuencia y garantía del mismo, etc. Sin embargo, esta última crítica no es tan sólida como podría parecer, ya que señala la dificultad de diseñar un pliego de condiciones detallado por parte del sector público, un problema al que también se tendría que hacer frente en caso de que el servicio fuera ofrecido por la iniciativa pública.

\paragraph{Vigencia de la concesión}
Sin embargo, las mayores dificultades se presentan en el campo del periodo de vigencia de la concesión. Supongamos que la subasta afecta a la prestación de un servicio que exige, como es frecuente, cuantiosas inversiones en activos específicos.

Si el periodo de vigencia es inferior a la amortización de los activos específicos se plantean problemas tanto si los activos son de propiedad privada como si son públicos. Si los activos son de titularidad pública, el regulador sólo tendrá que especificar las condiciones de operación de los activos, pero el concesionario no tendrá incentivos a utilizar eficientemente el equipo (conservación, tasa media de uso, etc.).\footnote{%
    Por ello, será imprescindible lograr que los gestores del equipo y el concesionario compartan el riesgo, pero la forma óptima de lograr esto puede ser enormemente compleja. Además, si los activos incorporan tecnologías muy especializadas, el primer concesionario tendrá ventajas sobre sus competidores en la segunda subasta por el aprendizaje realizado con dicho equipo capital, de forma que el grado de competencia ex post puede verse muy mermado.
} Por el contrario, si son de titularidad privada: (i) las empresas incorporarán en sus ofertas la cobertura del importante riesgo de que la concesión no sea renovada tras el primer período; (ii) además, la empresa no renovada tendrá incentivos a no vender su equipo en condiciones competitivas al nuevo concesionario, lo que implicará una ventaja no competitiva de la misma.  

Por el contrario, si el periodo es suficiente para garantizar la amortización de los activos específicos se suavizarían los problemas señalados, pero estos plazos suelen ser muy largos (por ejemplo, no menos de 20 años en generación de electricidad), con lo que las ventajas de la competencia potencial desaparecerían. Además, a medida que el plazo es más largo, las incertidumbres sobre las condiciones de demanda y costes (cambio técnico) aumenta, de modo que podría resultar imposible diseñar, poner en práctica y exigir el cumplimiento de un contrato óptimo a largo plazo. 

\clearpage
\section{Poder de Mercado: Relaciones verticales}

\textcolor{red}{\textbf{ESTE APARTADO SE MEJORA PERO SE INTRODUCE AQUÍ COMO COMODÍN. SU LOGAR ADECUADO SERÍA EL TEMA 15, COMO EXTENSIÓN TEÓRICA DE LA NTOI A LA TEORÍA DE LOS MERCADOS IMPUGNABLES Y LAS BARRERAS A LA ENTRADA}}

Muchas de los bienes y servicios con los que interactuamos diariamente no operan en mercados en los que exista un poder verdaderamente ostensible, ¿o sí? Paremonos a pensarlo. 

Pensemos en nuestra actividad diaria. ¿Cuántas tiempo dedicamos a nuestro teléfono? ¿Cuánto tiempo pasamos consumiendo entretenimiento en plataformas de streaming? ¿Cuánto tiempo pasamos realizando consultas en internet? Todas estas preguntas conducen a una clara conclusión: dedicamos mucho más tiempo del que creamos al consumo de bienes y servicios en los que interviene un alto grado de poder de mercado. Por suerte, regulado. Veámos porque.

En primer lugar, el número de proveedores/clientes/amigos con los que interactúamos a través de las TIC puede resultar enorme. Sin embargo, las plataformas a través de las cuales lo hacemos son un número asombrosamente reducido de empresas (en concreto, siete). En segundo lugar, aunque más rebuscado, en todas ellas subyace una tecnología característica de la Teoría tradicional (y, por ende, moderna) de la regulación: las telecomunicaciones. Por su condición de industria red, la industria de las telecomunicaciones presentas unas características muy concretas que exigen su regulación: (i) la estructura convexa de costes hace necesaria la condición de monopolio natural en lo que respecta a las redes de interconexión; y, (ii) tanto dentro como a través de estas redes de conexión interactúan/compiten un grado muy elevado de empresas en condiciones asimilables, en muchos casos, a la competencia perfecta, por lo que interactúan dos (o más mercados), uno con un poder de mercado enorme y otro en condiciones cercanas al Óptimo Paretiano.

Este ejemplo, que a continuación explotaremos, refleja dos lineas de importantes de investigación en el ámbito de la regulación: (i) los \textit{multi-sided markets} o plataformas; y, (ii) la liberalización de sectores, en particular, las convivencia entre diferentes mercados.

\subsection{Liberalización de sectores: Debates y retos}
¿Qué problemas pueden emerger cuando la empresa titular de las redes nacionales de comunicación vende servicios a diferentes empresas, por ejemplo operadoras de larga distancia, que operan en competencia y que a su vez venden al consumidor final?

En primer lugar, es obvio que los competidores en el segmento liberalizado necesitan acceso a los servicios que provee el monopolista natural, pero ¿y si este quiere operar en el segmento competitivo con su propia marca? ¿qué pasa si decide discriminar el acceso favoreciendo a su propia marca? O, por otro lado, imaginémos que la demanda de servicios crece de manera exponencial, ¿qué incentivos tiene el monopolista para mejorar y ampliar su infraestructura? En España, por ejemplo, el volumen de datos por linea móvil mensual crecio en casi un 200\% entre 2014 y 2021, ¿por qué iba el monopolista a ampliar la infraestructura para dar cabida a la creciente demanda?

Estos son solo algunos de los problemas que emergen en un escenario de estas características. Cada uno de ellos respresenta un gran núcleo o area donde la situación del monopolista constituye un debilidad enorme para el correcto funcionamiento del mercado. Por eso, en terminología de la NTOI, diremos que estas industrias \textbf{generan distorsiones/barreras verticales de acceso} (comportamientos que merman la competencia en el segmento liberalizado) y que el monopolista constituye un \textbf{upstream bottleneck}, un cuello de botella aguas arriba (ausencia de incentivos en mejorar la infraestructura, etc.).

En definitiva lo que queremos mostrar es que las industrias red y, más en concreto la liberalización de éstas, requiere analizar varias cuestiones que comprometen su correcto funcionamiento.

\subsubsection{Distorsiones verticales de acceso: Separación vertical}
Uno de los principales debates regulatorios que surgen es si debe limitarse la participación del monopolista incumbente en las actividades que se liberalizan y, accesoriamente, cómo garantizar que las condiciones de acceso al sector liberalizado sean las adecuadas para todos los participantes y garanticen la libre competencia.\footnote{%
    El argumento a favor de restringir la participación de los monopolistas regulados en otros mercados competitivos se basa en que estos pueden aprovecharse de los beneficios que obtienen en el mercado regulado para obtener una ventaja frente al resto de competidores (esto es, podrían aplicar subsidios cruzados) y extender así su poder de mercado del segmento regulado al competitivo. En realidad, esta cuestión también puede surgir en otros mercados en los que el monopolista no suministre inputs a sus competidores; basta con que los productos sean complementarios.
}

\paragraph{Segmento liberalizado: participación}
Evidentemente si se permitiera al monopolista competir en el segmento competitivo podría tratar de extender su poder de mercado mediante diversas estrategias, por ejemplo: (i) prácticas discriminatorias, aumentando el precio de acceso, reduciendo la calidad o, en caso extremo, negándose a suministrar acceso a su red a las empresas que compiten con él aguas abajo;\footnote{%
    Si los bienes no tuvieran una relación vertical, sino que sus consumos fueran complementarios. hablaríamos de vinculación o tying. Alternativamente, el monopolista podría emplear los beneficios del mercado regulado para reducir los precios en el mercado competitivo por debajo de sus costes, lo que se conoce como "precios predatorios". Todas estas prácticas supondrían infracciones del derecho de la competencia por abuso de posición dominante (en el caso de España, infringirían el artículo 2 de la Ley de Defensa de la Competencia). 

En última instancia, una vez reducida la competencia en el segmento competitivo, el monopolista aumentaría el precio del bien sujeto a competencia.
} y/o, (ii) reasignación de costes, en la medida en que algunos de los costes del monopolista sean comunes a la actividad competitiva y no competitiva, éste puede asignar artificialmente los costes a la no competitiva, lo que podría aumentar el precio regulado y proporcionar unos beneficios adicionales al monopolista.\footnote{%
    Si, además, el monopolista logra transferir parte de los costes variables al mercado regulado, se produce un subsidio cruzado que mejora su posición en el segmento competitivo, permitiéndole bajar precios y ganar cuota de mercado.
} Cualquiera de estas estrategias genera ineficiencias, distorsionando los precios de los mercados competitivo y no competitivo y reduciendo la cuota de mercado de las empresas competidoras. Si resulta que estas últimas son más eficientes, la ineficiencia será todavía mayor.\footnote{%
    En el sector de telecomunicaciones en EE. UU, se consideró que pesaban más estos efectos, dando lugar a la escisión de AT\&T en 1982 (las autoridades obligaron a que la empresa vendiera sus activos en el mercado no competitivo de llamadas locales y compitiera únicamente en el servicio liberalizado de llamadas de larga distancia).
}

No obstante, prohibir que las empresas reguladas operen en mercados con competencia supone también algunos costes en eficiencia que obliga a replantearse la prohibición: (i) si existen economías de escala en la producción de ambos bienes, entonces la prohibición impediría operar a la empresa más eficiente; (ii) la integración vertical del monopolista podría ser deseable como mecanismo para reducir los costes de transacción entre ambos segmentos; y, (iii) si la competencia en el segmento liberalizado es imperfecta, entonces permitir la entrada del monopolista sería socialmente deseable, al aumentar la producción y disminuir los precios.

\paragraph{Segmento liberalizado: acceso}
En lo relativo al acceso al segmento liberalizado del mercado, el regulador debe velar porque se produzca de forma no discriminatoria. Con esto en mente, existen diferentes alternativas para la regulación del precio óptimo de entrada. Entre las más utilizadas están:
\begin{la}
    \item \textbf{Efficient price-component rule}. Inicialmente propuesta por BAUMOL y WILLIG, trata de evitar que la entrada de nuevos competidores genere ineficiencias productivas. Consiste en dos sencillos pasos que deben ser aplicados por el regulador: (i) el precio de acceso a la red debe ser igual al coste de oportunidad que le supone al monopolista conceder el acceso; y, (ii) el coste de oportunidad está formado por la suma del coste marginal del servicio y el coste de oportunidad que le supone perder ventas en el mercado competitivo frente al entrante (la pérdida marginal de demanda). Una vez se dispone de estos datos, se puede fijar el precio de acceso;
    \item \textbf{Precios de acceso óptimo}.\footnote{%
        Para un desarrollo más detallado del modelo y de posibles extensiones se puede consultar CHURCH y WARE (2000). p. 873.
    } Por otro lado, ARMSTRONG, DOYLE, and YICKERS (1996) propusieron un mecanismo de fijación alternativo que tuviera en cuenta el precio del activo esencial (como en ECPR), pero también el precio del mercado en competencia. Con el fin de garantizar la viabilidad del monopolista, el regulador fijará: (i) un precio de acceso igual al coste marginal como óptimo de primer orden, si garantiza la viabilidad; y, (ii) un precio óptimo de segundo orden atendiendo a la regla RAMSEY-BOITEAUX si existen economías de escala y/o de alcance o la viabilidad del monopolista no queda garantizada de la otra forma.
\end{la}

En cualquier caso, el precio de acceso óptimo es inferior al que fijaría el monopolista en ausencia de regulación: sin un regulador que fije estos peajes, las empresas tenderán a escoger un peaje demasiado alto desde el punto de vista de los consumidores. 

\subsection{Multi-sided Platforms: Economía digital}
Por otro lado, y relativo a la otra linea de investigación que hemos introducido al principio de la sección, la creciente importancia de las plataformas digitales en nuestras economías ha dado lugar a que cada vez haya más interés en estudiarlas desde una perspectiva tanto teórica como aplicada (política de competencia y regulación).\footnote{%
    La modelización de las plataformas tiene su origen en TIROLE y ROCHET (2003), donde exponen la noción de \textit{two-sided markets}, esto es, una plataforma que pondría en común a dos grupos de agentes.

Si bien la teoría se ha generalizado, extendiéndose al concepto de multi-sided markets, platforms o matchmakers (EVANS y SCHMALENSEE, 2016). En el contexto actual de crecimiento de la economía digital, la mayoría de las Big Techs cumplen con las características de una plataforma, si bien el concepto también podría aplicarse a situaciones no digitales. Puede pensarse en una discoteca, una tarjeta de crédito, la moneda de curso legal de un país o una cadena de televisión. En aras de una mayor sistematicidad, se podría plantear una clasificación de distintos tipos de plataformas que no busca ser completa ni general:
\begin{la}
    \item \textbf{Facilitadoras de intercambios}. Determinadas plataformas tienen como principal objetivo poner en contacto a dos agentes para que realicen un intercambio (Wallapop, Tinder, BlaBlaCar, etc);\footnote{El caso de Amazon es representativo de las vertiginosas dinámicas de los mercados digitales donde, en un primer momento, nos encontraríamos con un facilitador de intercambios de libros y, en la actualidad, con un vendedor de productos (donde ya no se darían con tanta intensidad las características puras de plataforma).
    }
    \item \textbf{Financiadas por publicidad}. Son plataformas en las que un grupo de agentes está interesado en mostrar publicidad a otro grupo de agentes y, como veremos a continuación, pagan por ello (Facebook, Telecinco, 20 Minutos); 
    \item \textbf{Sistemas de transacciones}. Permiten realizar transacciones entre distintos grupos de usuarios (Visa, el euro, un criptoactivo); y,
    \item \textbf{Plataformas de software}. Ponen en común a desarrolladores de programas informáticos o Apps con los potenciales usuarios (el sistema operativo Windows, iOS).
\end{la}
}


Una plataforma es un tipo de empresa cuya función es poner en común a varios grupos de usuarios con distintas características. En un contexto de altos costes de transacción, donde a los distintos usuarios les es difícil o costoso encontrar una contraparte para realizar un intercambio económico, la plataforma les puede ayudar a encontrarlo a cambio, o no, de una contraprestación.\footnote{%
    Según el teorema de Coase, en ausencia de fricciones y costes de transacción, el resultado de la negociación entre un oferente y un demandante será óptimo. No obstante, en presencia de altos costes de búsqueda, no se produce un emparejamiento directo del oferente y del demandante, por lo que surge un nuevo actor, la platafonna intennediaria, la cual ayuda a las partes a reducir los costes de transacción. Por ejemplo, la aplicación The Fork reduce los costes de búsqueda para los comensales y reduce los costes de gestión de las reservas para los restaurantes.
}

La caracterización económica de las plataformas difiere de las empresas neoclásicas en varios puntos.

\subsubsection{Características empresariales: Barreras de entrada}
En primer lugar, y relativos a las características intrínsecas de su actividad empresarial, las plataformas se ceracterizan por la existencia de los siguientes efectos:
\begin{lnum}
    \item \textbf{Network effects}. En la mayoría de los casos, el valor añadido que aporta la plataforma es satisfacer el deseo de un grupo A de interactuar con otro grupo B. En la medida en que dicho grupo B es más numeroso, el valor de la plataforma aumenta para el grupo A (por ejemplo, un periódico de tirada nacional en vez de provincial, siendo el grupo B los lectores y el grupo A los anunciantes en el periódico). No obstante, los efectos de red también podrían ser internos al grupo B (los usuarios están en lnstagram debido a que sus amigos también lo están, independientemente del lado A de anunciantes e instagramers) o incluso recíprocos entre el lado A y el lado B (en Wallapop, ambos grupos se benefician de que haya tanto muchos compradores como muchos vendedores). En este sentido, BELLETLAMME (2021) diferencia entre cross-group y with-in group network effects. Para el autor, la gestión de los efectos de red de los distintos grupos de usuarios es la actividad característica de una plataforma;
    \item \textbf{Switching-costs}. Elevados costes de cambiar de empresa (switching costs). En algunos casos, una vez que el usuario ha empezado a interactuar con una determinada plataforma, le puede ser muy costoso cambiar de plataforma (en términos de tiempo, costes de búsqueda, aprendizaje, etc). Además, los datos personales ya acumulados por una plataforma juegan un papel importante a la hora de retener a sus clientes (por ejemplo, el hecho de tener listas de reproducción creadas en Spotify);
    \item \textbf{Critical mass: Economías de escala}. En especial en el caso de las plataformas digitales, el coste marginal de añadir un usuario más a la provisión del servicio es cercano a cero. Además, un gran número de plataformas presentan el conocido como problema de la masa crítica, esto es, para que la plataforma pueda crear valor, es necesario que reúna una base significativa de clientes en todos los grupos relevantes. Por lo tanto, deben recurrir a una serie de estrategias para alcanzar dicha masa crítica. Las plataformas financiadas mediante publicidad tienden a realizar una actividad de generación de contenido para los usuarios previa a la monetización de los espacios publicitarios. Otra solución podría ser la estrategia de dividir y vencer (JULLIEN, 2011), según la cual la plataforma subsidiaría a los agentes que pertenezcan al grupo con una mayor elasticidad-precio, para que se sumen más fácilmente a la plataforma y, así, una vez la plataforma haya ganado escala de clientes, poder atraer a otros grupos;\footnote{%
        Por ejemplo, Diners Club, la primera empresa emisora de tarjetas de pago en los años 50's, se vio enfrentada al problema de cómo conseguir que los consumidores adoptasen el pago por tarjeta. Inicialmente comenzó distribuyendo tarjetas para los comensales de una serie de restaurantes en Nueva York para que pagasen sus cenas y así se acostumbrasen a esta nueva forma de pago. Una vez se construyó una base sólida de clientes, el negocio se expandió, convirtiéndose en una tarjeta de pagos aceptada en la mayoría de tiendas.
    }
    \item \textbf{Big data y algoritmos}. Además de las barreras de entrada ya mencionadas (efectos de red, economías de escala, switching costs) es importan te hacer una referencia al papel de los datos como fuente de ventaja competitiva en los mercados digitales. Como apuntan JONES y TONETTI en un influyente artículo (2020),\footnote{\href{}{\textit{Nonrivalry and the Economics of Data}. JONES y TONETTI (2020)}
    } los datos presentan características de no rivalidad, esto es, pueden ser utilizados por distintos agentes económicos sin que su dotación disminuya. En otras palabras, el análisis masivo de datos presenta características de bien público. Por lo tanto, hay incentivos por parte de las empresas a acumular datos y a evitar que otras empresas accedan a estos, generando así barreras de entrada significativas, tanto en mercados digitales como tradicionales. Una empresa con amplios datos puede servir mejor las necesidades de los clientes, mejorar sus procesos internos, extraer un mayor excedente del consumidor, etc. Una empresa entrante en un mercado puede tener grandes dificultades para competir sin estos conocimientos que arrojan los datos (\textit{data is the new oil}). Los autores apuntan a que los rendimientos crecientes a escala que presentan los datos permitirían un mayor bienestar social en caso de que los datos fuesen considerados un cuello competitivo (como en las industrias de red tradicionales) y su acceso fuese permitido a todas las empresas que compiten en un mercado. Como veremos, estas ideas han influenciado en gran medida a las nuevas regulaciones en esta materia.  
\end{lnum}


Esto tiene una implicación fundamental sobre la determinación del precio óptimo: ya no dependerá solo de la elasticidad de la demanda y los costes marginales de cada lado del mercado, sino que habrá que tener en cuenta cómo responderán las demandas del resto de lados.\footnote{%
    Puede incluso ser óptimo fijar un precio nulo o negativo en un lado del mercado para impulsar la demanda y los precios del otro lado.  Por ejemplo, una plataforma online de venta de productos, al fijar las comisiones que cobrará a los vendedores y los precios que pagarán los compradores, debe tener en cuenta que los compradores valoran que haya variedad de empresas y, a su vez, los vendedores valoran que haya muchos potenciales compradores. Por tanto, aumentar el precio a un lado del mercado no solo reducirá la demanda de ese lado, sino que dará lugar a una caída de la demanda del otro lado y así sucesivamente.
} 

\subsubsection{Características de mercado: }


Por ello, los mercados de varias caras suelen tener una tendencia a la concentración, pudiendo incluso convertirse en monopolios debido a la interrelación de varios efectos:
	Los efectos de red indirectos pueden generar círculos virtuosos de demanda.
	Estrategias de las plataformas para hacer cautivos a sus usuarios y excluir así a los competidores como dificultar su compatibilidad, la publicidad, la innovación o adquisición de empresas innovadoras, la integración vertical y discriminación para favorecer sus propios productos, etc.


Aunque en la mayoría de los mercados, estos efectos sobre la competencia se tratarían con la política de competencia, la rapidez a la que evolucionan ha llevado a muchos países a optar por un enfoque regulatorio para resolver estos problemas. La iniciativa más ambiciosa en los últimos años es el Reglamento de Mercados Digitales (Digital Markets Act, DMA), que entró en vigor a partir de 2023 sobre las grandes plataformas digitales que operen en la UE.  
La DMA impone una serie de obligaciones especiales dirigidas a asegurar que las grandes plataformas digitales no puedan imponer condiciones no equitativas, para lo que tendrán que desarrollar proactivamente ciertas prácticas (regulación) y abstenerse de ciertas conductas (relaciones verticales). Entre las medidas, se pueden destacar prohibiciones de cláusulas de paridad, de vinculación, de auto-preferencia, de aplicar condiciones desventajosas y obligaciones de transparencia, de portabilidad de datos o de interoperabilidad.

5.4.1. Aplicaciones
La Comisión Europea, en aplicación de la Ley de Mercados Digitales (Reglamento 2022/1925, más conocido por sus siglas DMA), ha sancionado e impuesto una serie de obligaciones para preservar la competencia en los citados mercados, a:
	APPLE. Con una multa de 500 millones de euros, por las restricciones impuestas sobre los usuarios de su red (aguas abajo), que restringen la competencia al uso exclusivo de su propia red, con las ventajas que esto supone, tanto a nivel de ingresos como a nivel de control;  y,
	META. Con una multa de 200 millones de euros. Por utilización de su poder de red (más que el poder de mercado, por lo tanto estamos ante la otra cara de la plataforma, i.e. consumidores) para restringir las condiciones de consumo posible artificialmente (bundling) sin alternativa.  En virtud de la DMA, los guardianes de acceso (gatekeepers) deben solicitar el consentimiento de los usuarios para combinar sus datos personales entre servicios. Los usuarios que no den su consentimiento deben tener acceso a una alternativa menos personalizada, pero equivalente..
Las sanciones son las primeras bajo la DMA, la nueva norma de la Unión Europea para regular el poder de mercado de las grandes plataformas de internet.





 












\clearpage
\section{Conclusión} 

\subsection{Extensiones y relación con otras partes del Temario}


\subsection{Opinión}



\subsubsection{Idea final (Salida o cierra)}


\clearpage
\pagenumbering{gobble}
\MyBibliography
\MyBibItem{Autor}{Título}{Año}{DOI -- LINK}




%ANEXOS
\clearpage
\anexo{\textbf{Preguntas Test}}
%\input{\TemaCode/questions.tex} 



\clearpage
\anexo{Regulación: Información perfecta}\label{anx:informacion_perfecta}
3. Regulación con información perfecta
Puntos clave:
	Las teorías tradicionales se basan en las justificaciones normativas de la regulación, con:
	Una empresa regulada que trata de minimizar sus costes; y,
	Un regulador que trata de maximizar el Bienestar Social 
Sin que existan asimetrías informativas ni comportamientos estratégicos.
	Entre estas teorías estarían las que buscan que los precios regulados sean óptimos de primer orden: precios según el coste marginal, precios no lineales o precios Pico-Valle. 
En estos casos, no existiría un trade-off entre alcanzar la solución eficiente en términos de precios y cubrir los costes del monopolista.
	Si surge el mencionado trade-off, habrá que aplicar el teorema del second best con el fin de alcanzar unos precios óptimos de segundo orden que serán superiores al coste marginal y por tanto, generan ineficiencias. La regla de precios más conocida en esta línea son los precios de RAMSEY-BOITEUX.
	Las teorías tradicionales han sido la base teórica para esquemas regulatorios como la regulación mediante el coste del servicio (cost plus) todavía ampliamente utilizados.
La literatura tradicional suele partir de:
	La empresa regulada tiene características de monopolio natural;
	 Suministra uno o más servicios y minimiza sus costes dados unos niveles de tecnología, producción y unos precios de los inputs (no existe ineficiencia X).
	Un regulador, plenamente informado sobre los costes y la demanda de la empresa regulada, cuyo objetivo es fijar los precios

2.2. Simple two-parts tariffs and non-linear prices
Otra alternativa que puede permitir alcanzar un resultado eficiente es relajar la necesidad de que la empresa regulada cobre un precio lineal por cada unidad vendida. En este caso, el regulador deseará encontrar el esquema tarifario o los diferentes niveles de precios de manera que se maximice el beneficio social sujeto a la restricción de participación del monopolista. 
2.2.1. Modelo: Simple Two-parts tariffs 
El caso más sencillo es la tarifa en dos partes, según la cual el pago del consumidor por la compra del bien consta de dos componentes:
	Un pago fijo, que cabe identificar con una cuota de entrada en el mercado, y 
	Un pago unitario constante por la cantidad de producto demandada.
2.2.1.1. Supuestos
Suponemos que:
	Hay N consumidores idénticos, cada uno con una demanda una cantidad, q_i,  y cuenta con un excedente bruto S_i; y, 
	La empresa regulada tiene un coste marginal c y unos costes fijos, de modo que su función de coste total es C= F+ c·q. 
2.2.1.2. Desarrollo analítico
En este contexto, una tarifa en dos partes que permitiría alcanzar un resultado eficiente y cubrir los costes totales del monopolista sería fijar una cuota de acceso a cada consumidor A= F/N y un pago unitario p = c, la siguiente tarifa para el consumidor i:
τ_i=A+pq_i=F_0/N+cq_i
Con este esquema de precios  se alcanza una solución óptima de primer orden donde: (i) en términos marginales, cada consumidor paga un precio de uso igual al coste marginal del monopolista; y, (ii) el pago fijo actúa como un impuesto de suma fija que permite cubrir la brecha entre los pagos unitarios y los costes totales.
2.2.1.3. Implicaciones: Exclusivistas o no exclusivistas
El bien solamente se proveerá si A< (S(q_i )-p·q_i,): si el excedente neto de cada consumidor es mayor que el pago fijo. De lo contrario, no tendrá sentido económico hacerlo puesto que el coste asociado sería mayor que el excedente bruto de los consumidores.  Por tanto, en presencia de varios tipos de consumidores, la imposición de un pago fijo excluirá del mercado a las clases de consumidores con menor disponibilidad a pagar, dado que el excedente que genera su participación en el mercado es inferior a la cuota fija. Si suponemos que el monopolista debe cobrar a todos los consumidores la misma tarifa y el mismo precio. Tendremos diferentes consecuencias en función de la tarifa de acceso establecida por el monopolista.
Si el monopolista decide cobrar a todos los consumidores la tarifa correspondiente a la de los clientes de menor demanda y un precio igual al coste marginal, el monopolista no se apropia completamente del excedente del consumidor:
 
Si el monopolista le compensa establecer una tarifa de entrada elevada (correspondiente al excedente del consumidor de los consumidores de demanda alta), dejará fuera del mercado a los consumidores de demanda baja (exclusivista):
 
2.2.2. Modelo baseline: non-linear pricing
La solución al problema de abandono de algún tipo de consumidor que se acaba de señalar viene dada por el ofrecimiento a los consumidores de las siguientes alternativas tarifarias: 
	Un precio uniforme superior al coste marginal p_1=(c + t); o,
	Una cuota fija (A) más un precio unitario igual a dicho coste marginal p_2=c. 
Quienes consuman por encima de una cantidad determinada (siendo esta: A = t·q) preferirán la segunda opción, y los restantes la primera, por tanto, este menú tarifario es una mejora paretiana respecto a la situación de precio uniforme y también respecto a la tarifa en dos partes. Sin embargo, a pesar de esta simplificación, los esquemas de precios no uniformes pueden ser mucho más complejos. Las tarifas podrían constar de múltiples componentes, y la opción entre distintas alternativas tarifarías puede dar lugar a tarifas por bloques. 
2.2.3. Modelo TIROLE: optimal non-linear pricing
A continuación, se muestra un breve desarrollo analítico del modelo de tarifas totalmente no lineales, siguiendo a TIROLE (1988), que muestra los resultados principales que se alcanzan.
2.2.1.1. Supuestos
Partimos de los siguientes supuestos:
	El precio depende de la cantidad consumida, T(q). Siendo T el precio total pagado por la cantidad q;
	Existen dos grupos de consumidores, a (alta valoración) y b (baja), con utilidades: [θ_a·U(q)-T] y [θ_b·U(q)-T], respectivamente, con θ_a>θ_b;
	Una masa de consumidores igual a 1, sin que el monopolista pueda identificar a qué grupo pertenece cada consumidor, existiendo una proporción Θ_a y Θ_b de cada grupo.
2.2.1.2. Desarrollo analítico
Con tarifas plenámente no lineales podemos elegir cualquier tipo de función para T(q), con tal de que: 
	(q_a,T_a ) sea consumido por el consumidor a; y
	(q_b,T_b )  sea consumido por e1 consumidor b
Para resolver dicho programa debemos saber qué restricciones se cumplen en forma de igualdad. 
2.2.1.3. Desarrollo gráfico
  
 
2.2.1.4. Implicaciones
En equilibrio, la tarifa no lineal óptima conduce a descuentos por cantidad y asigna diferentes cantidades y pagos a cada tipo de consumidor:  
	El grupo con alta valoración (a) consume su cantidad eficiente (no distorsion at the top), si bien la tarifa no consigue extraer todo el excedente de este grupo, con lo que dicho grupo obtiene una renta positiva (excedente) derivada de su ventaja informativa; mientras que,
	El grupo con baja valoración (b) consume la cantidad q_b, inferior a la óptima (distorsion at the bottom), si bien T_b=θ_b·U(q_b), con lo que la tarifa consigue extraer todo el excedente de los consumidores con baja valoración. 
2.3. Peak-Load Pricing 
Muchas industrias reguladas emplean un mismo activo fijo (por ejemplo, plantas de generación eléctrica, líneas de transporte y distribución de electricidad, tuberías de gas, redes telefónicas o ferroviarias) para producir el mismo bien en distintos momentos del día o del año. Esto plantea una nueva problemática al regulador: los precios deben reflejar que los costes marginales de producción no son iguales dependiendo de si la capacidad está o no saturada.
2.3.1. Modelo
Vamos a plantear un modelo sencillo aplicado a este sector para ilustrar esto. Diremos que existe una configuración especial de la demanda que permite dividir el período total de producción en:
	Un subperiodo en que la demanda es alta, o peak. Por ejemlpo, la demanda eléctrica durante las horas de día, q_D (p_D); y,
	Otro subperiodo en que es baja, o load: la demanda eléctrica durante las horas de noche, q_N (p_N).
2.3.1.1. Supuestos
Suponemos que:
	Para cualquier precio la demanda es mayor durante el periodo pico (horas de día) que durante el valle (horas de noche).
	El bien no es almacenable, de forma que los consumidores no pueden comprar en un subperiodo para consumir en otro (no es posible tampoco la reventa entre consumidores). 
	En cada periodo la empresa determina la capacidad de producción K (suponemos que es divisible), con un coste de inversión C_K K, y produce electricidad con un coste operativo marginal C_E. 
	El excedente bruto en cada periodo (área bajo la curva de demanda) viene dado por S(q_i), y tenemos que S^' (q_i )=p_i;
2.3.1.2. Desarrollo analítico
Los precios óptimos se obtienen solucionando un problema de maximización del excedente social (〖EC_"bruto" -C〗_K K) sujeto a que la producción en cada periodo debe ser menor o igual a la capacidad instalada. El desarrollo analítico del problema completo se puede ver en [Ver Anexo III]. 
2.3.2. Implicaciones
2.3.1. Primer caso: resultados típicos del modelo pico-valle
Un resultado típico (y válido) de este esquema tarifario divido en periodos conllevaría que:
	Durante el periodo valle, la demanda es menor a la capacidad y el precio solo cubre los costes marginales operativos y no aporta nada a los de capital.  
	Mientras que en el pico, los precios permiten cubrir tanto los costes marginales como los costes de capital, desplazando el precio durante las horas de demanda alta (pico).
 
2.3.2. Segundo caso: deplazamiento del pico
Si las demandas en ambos periodos son tan elásticas que al aplicar la regla de precios del modelo tradicional Pico-Valle todos los costes de capital son soportados por los consumidores del periodo pico, la demanda pico cae y la demanda valle aumenta de forma que los picos y los valles se invierten. 
En este caso, la regla de precios óptima implica que los costes marginales operativos y de capital (capacidad) se reparten entre los consumidores de los periodos pico y valle.  
 
2.3.3. Valoración y extensiones
La literatura  ha incluido varias complicaciones para dotar de mayor realismo a estos modelos:
	Posibilidad de elección (en el caso de la electricidad) entre distintas tecnologías de generación con distintas ratios capital-combustible.
	La demanda no se divide únicamente entre peak y load, sino que varía continuamente entre un máxímo y un mínímo. 
	Introducción de aleatoriedad en demanda y oferta (por ejemplo, con apagones imprevistos o necesidades de mantenimiento).


{Rate of return}
Una de las primeras aportaciones en esta linea fue el modelo de AVERCH-JOHNSON (1962). Si asumimos que el regulador no tiene ninguna función objetivo y que solo conoce el coste de capital de la empresa ($r$) y, por tanto, carece de información sobre la función de producción de la empresa, sus costes o su demanda y considerando que los mercados debían regularse a través del precio, y cómo ampliación lógica de los modelos coste del servicio,\footnote{%
    La regulación basada en el coste del servicio ha sido la tradicionalmente aplicada en EE.UU. y otros países occidentales durante buena parte del siglo XX. Una referencia excelente para ampliar los detalles de este tipo de regulación es CHURCH y WARE (2000)

Los modelos \textit{cost of service} tienen su razón de ser en la fijación de precios. Si regulamos un mercado determinando el precio al que deben ofrecer el producto los ofertantes, ¿cómo establecer dicho precio sin que los ofertantes cierren su persiana? La idea consistía en que el precio debía compensar a las empresas el coste incurrido al prestar el servicio:
\begin{lnum}
    \item \textbf{Determinar el coste del servicio}. Resultado de sumar los costes operativos (costes administrativos, laborales, impositivos, de materias primas, etc.) y los costes de capital (depende del stock de capital invertido, la tasa de rendimiento permitida y la depreciación);
    \item \textbf{Fijar un precio de prestación}. Se fija un precio que permite cubrir los costes totales CT(Q), de modo que P·Q = CT(Q).
\end{lnum}
El precio estipulado, bien sea al inicio o durante el transcurso del tiempo, acaba siendo igual a los costes operativos por lo que en realidad es un esquema de revelación de la información privada. En la actualidad este esquema se sigue usando en el ámbito de la provisión de servicios por parte de empresas públicas.
} el modelo de AVERCH-JOHNSON gira en torno a la tasa de retorno óptima del capital (rendimiento óptimo) que debe fijar el regulador con el fin de preservar/incentivar la prestación sin generar grandes distorsiones al proveedor.

{Desarrollo analítico}
Suponiendo que: (i) el monopolista produce bajo una tecnología de producción que utiliza únicamente trabajadores y capital; (ii) que persigue maximizar sus beneficios escogiendo la cantidad óptima de factores productivos; y, (iii) que el regulador puede fija una tasa máxima de rendimiento del capital, $s$, mayor que el coste de oportunidad del capital, $r$, pero menor que la tasa de retorno que percibiría un monopolista sin regular, $r_m$. Podemos formular el problema de maximización del monopolista como sigue:


$$\begin{aligned}
    \boxed{\begin{aligned}
        &\max \Pi(K,L)=P_x^d(x)F(K,L)-wL-rK\\[0.5em]
        &\text{s.a.}\quad \Pi(K,L)\ge 0
    \end{aligned}}
\end{aligned}$$


{Implicaciones de Política Económica}
Desde una perspectiva de Política Económica, la regulación del precio mediante cobertura (bien sea de costes o de retorno del capital) genera una serie de problemas razonables:
\begin{lnum}
    \item \textbf{Ineficiencia X}. En primer lugar, no se genera incentivos a la reducción de costes, ya que ésta mejora no repercute en beneficio del regulado. De hecho, se pueden crear incentivos para expandir los costes de forma beneficiosa para los gestores de la empresa regulada, dado que éstos saben que todos los costes justificados por la empresa serán cubiertos por los precios;\footnote{%
        Además, se está regulando como beneficio la remuneración del capital, que como sabemos no equivale al concepto de beneficio que se emplea en la ciencia económica. Por tanto, bajo esta práctica regulatoria, la reducción de costes implica que el precio regulado se verá a su vez reducido, dejando a la empresa con la misma tasa de beneficio. Esto puede incentivar la generación de gastos inútiles para no superar el rendimiento máximo.
    }
    \item Tampoco se incentiva la innovación en términos de nuevos productos o servicios. Ya que la empresa no percibirá, en su caso, los mayores beneficios asociados;
    \item \textbf{Precios RAMSEY-BOITEAUX}. Por último, como no están basados en los costes marginales ni reflejan las elasticidades de la demanda, estos mecanismos no garantizan que se alcancen los precios RAMSEY-BOITEUX que constituyen el second-best teórico;
\end{lnum}

Pese a sus deficiencias, ésta es una técnica habitualmente utilizada en la regulación de servicios públicos,\footnote{%
    En España, la mayoría de los monopolios naturales regulados se retribuyen siguiendo un esquema de coste del servicio. Por ejemplo, las tarifas de los servicios prestados por ADIF relativos al suministro de corriente de tracción a trenes o la carga y descarga de contenedores de mercancías, están sometidas a esta regulación de precios. Lo mismo sucede con la retribución del transporte eléctrico. Véase, por ejemplo: Ministerio de Hacienda, Coste efectivo de los servicios prestados por las entidades locales. \href{https://www.hacienda.gob.es/es-ES/CDI/Paginas/InformacionPresupuestaria/InformacionEELLs/CosteServicios.aspx}{Coste de los servicios públicos locales. Ministerio de Hacienda}
} posiblemente por sus escasos requerimientos informativos y (relativamente) fácil comprobación a posteriori mediante auditoría, ya que sólo requiere conocer beneficios y stock de capital.

{Efecto AVERCH-JOHNSON}
Además, el resultado fundamental es que se produce el conocido como efecto AVERCH-JOHNSON: los factores productivos se utilizan en proporciones que no minimizan los costes de producción (distorsión en la elección de los niveles de inputs) y, en particular, la empresa se sobrecapitaliza.\footnote{%
    Dado que se restringe el beneficio obtenible, ligándolo a la cantidad de capital utilizado, es rentable para el monopolista utilizar en cuantía excesiva el capital, como forma de aumentar la base sobre la que se pueden obtener ganancias.
En particular, la restricción de la tasa de beneficio abre una brecha entre el coste de capital real de la empresa y su coste en términos netos, después de considerar los beneficios de aumentar el stock de capital, que en términos marginales tiene un rendimiento positivo de (s-r).
}

Por ello, no solo no es sencillo saber qué tasa se elige en el intervalo mencionado sino que, además, se puede demostrar que no es seguro que la reducción de $s$ disminuya el grado de sobrecapitalización (aproximándola a la tasa de beneficio que se obtendría en condiciones competitivas).

\end{document}